\documentclass[12pt]{scrreport}
\usepackage[utf8]{inputenc}
\usepackage{graphicx}
\usepackage[a4paper, margin=2.5cm]{geometry}
\usepackage{hyperref}
\usepackage{amsmath}
\usepackage{enumitem}
\usepackage{mdframed}
\usepackage{amsthm}
\usepackage{amssymb}
\usepackage{mathtools}
\usepackage{cite}
\graphicspath{ {images/} }
\usepackage{xcolor}         % Extended colors
\usepackage{color}         % Color extended names
\usepackage{nomencl}
\usepackage{lipsum}
\makenomenclature
\renewcommand{\nomname}{Notation Index}



\newsavebox{\selvestebox}
\newenvironment{takeaways}
  {
   \begin{lrbox}{\selvestebox}%
   \begin{minipage}{15cm}
     \textbf{Takeaways}
   }
  {\end{minipage}\end{lrbox}%
   \begin{center}
\setlength\fboxsep{.5cm}
   \colorbox[HTML]{F8E0E0}{\usebox{\selvestebox}}
   \end{center}}
% Theorem handle


\newtheorem{theorem}{Theorem}[chapter]
\newenvironment{ftheorem}
  {\begin{mdframed}\begin{theorem}}
  {\end{theorem}\end{mdframed}}
\newtheorem{assumption}{Assumption}[chapter]
\newtheorem{corollary}{Corollary}[theorem]
\newtheorem{definition}{Definition}[chapter]
\theoremstyle{remark}
\newtheorem{example/}{Example}[chapter]
\newenvironment{example}
  {\renewcommand{\qedsymbol}{$\diamondsuit$}%
   \pushQED{\qed}\begin{example/}}
  {\popQED\end{example/}}
\theoremstyle{plain}
\newtheorem{lemma}{Lemma}[chapter]
\newtheorem{problem}{Problem}[chapter]
\newtheorem{proposition}{Proposition}[chapter]
\newtheorem{remark/}{Remark}[chapter]
\newenvironment{remark}
  {\renewcommand{\qedsymbol}{$\diamondsuit$}%
   \pushQED{\qed}\begin{remark/}}
  {\popQED\end{remark/}}
\newtheorem{subassumption}{}[assumption]

\renewcommand{\proofname}{\textbf{Proof}}
% definition of constants

\newcommand{\CP}{C_\P}
\newcommand{\Ctau}{C_\tau}
\newcommand{\LearnRate}{\varepsilon_n}

%
\DeclarePairedDelimiterX{\inner}[2]{\langle}{\rangle}{#1, #2}
\newcommand{\C}{\mathbb{C}}
\newcommand{\G}{\mathbb{G}}
\newcommand{\E}{\mathbf{E}}
\renewcommand{\P}{\mathbf{P}}
\newcommand{\R}{\mathbb{R}}
\newcommand\norm[1]{\left\lVert#1\right\rVert}

\title{
  {
    Robust Weighting and Matching Techniques for Causal Inference in Observational Studies with Continuous Treatment
  }
  \\
  {\large Universität Stuttgart}
  \\
  {\includegraphics{unistuttgart_logo_deutsch.jpg}}
}

\author{Ioan Scheffel}
\date{\today}

\begin{document}


\maketitle

\tableofcontents 

%\chapter{Introduction}
%% SAY THAT RESEARCHERS ARE OFTEN LEFT WITH OBSERVATIONAL STUDYS TO ANSWER THEIR QUESTIONS

% SAY SOMETHING ABOUT PS METHODS AS TO DO WITH CAUSAL INFERENCE IN OS

% SAY THAT MODEL DEPENDENCY IS OFTEN A PROBLEM WITH PS

% MENTION DIFFERENT APPROACHES TO ALLEVIATE THIS PROBLEM

% HIGHLIGHT THE METHODS BEING DISCUSSED IN THE THESIS

% OUTLINE STRUCTURE 
%% - WEIGHTING APPROACH BINARY
%%  - MODIFICATIONS IN BINARY SETTING TO ESTIMATE ATE
%%  - EXTENSION TO CONTINUOUS TREATMENT
%% - MATCHING APPROACH BINARY 
%%  - POSSIBLE EXTENSION
%% - SIMULATION STUDY
%% - REAL DATA USECASES


Researchers are often left with observational studies to answer questions about causality. When confounders are present the task of infering causality can become arbitrarily complex. Propensity score methods \cite{Rosenbaum1983}, e.g. inverse probability weighting or matching, are popular methods to adjust for confounders. Usually these methods rely heavily on estimates of the true propensity score, which are known to suffer from model dependencies and misspecification\cite{Kang2007}. This issue becomes more pressing when moving from binary to continuous treatment\cite{Hirano2005}. Therefore methods have been developed to directly target imbalances in the data\cite{Fong2018}\cite{Hainmueller2012}\cite{Zubizarreta2015}.
We take a closer look at \cite{Wang2019} and extend the analysis to settings with continuous treatment 
\cite{Vegetabile2020}\cite{Tubbicke2020}.


\chapter{Causal Inference}
In this chapter we want to give a introduction to causal inference.
We particularly highlight the role of propensity score analysis and explain its importance in observational studies.


\section{The Rubin Causal Model}
from wiki:
The Rubin causal model (RCM), also known as the Neyman–Rubin causal model,[1] is an approach to the statistical analysis of cause and effect based on the framework of potential outcomes, named after Donald Rubin. The name "Rubin causal model" was first coined by Paul W. Holland.[2] The potential outcomes framework was first proposed by Jerzy Neyman in his 1923 Master's thesis,[3] though he discussed it only in the context of completely randomized experiments.[4] Rubin extended it into a general framework for thinking about causation in both observational and experimental studies.[1]
\section{Propensity Score Analysis}
from wiki:
In the statistical analysis of observational data, propensity score matching (PSM) is a statistical matching technique that attempts to estimate the effect of a treatment, policy, or other intervention by accounting for the covariates that predict receiving the treatment. PSM attempts to reduce the bias due to confounding variables that could be found in an estimate of the treatment effect obtained from simply comparing outcomes among units that received the treatment versus those that did not. Paul R. Rosenbaum and Donald Rubin introduced the technique in 1983.[1]

The possibility of bias arises because a difference in the treatment outcome (such as the average treatment effect) between treated and untreated groups may be caused by a factor that predicts treatment rather than the treatment itself. In randomized experiments, the randomization enables unbiased estimation of treatment effects; for each covariate, randomization implies that treatment-groups will be balanced on average, by the law of large numbers. Unfortunately, for observational studies, the assignment of treatments to research subjects is typically not random. Matching attempts to reduce the treatment assignment bias, and mimic randomization, by creating a sample of units that received the treatment that is comparable on all observed covariates to a sample of units that did not receive the treatment.

For example, one may be interested to know the consequences of smoking. An observational study is required since it is unethical to randomly assign people to the treatment 'smoking.' The treatment effect estimated by simply comparing those who smoked to those who did not smoke would be biased by any factors that predict smoking (e.g.: gender and age). PSM attempts to control for these biases by making the groups receiving treatment and not-treatment comparable with respect to the control variables. 

from a paper:
Propensity score weighting is one of the techniques used in controlling for selection biases in non-
experimental studies. Propensity scores can be used as weights to account for selection assignment
differences between treatment and comparison groups. One of the advantages of this approach is
that all the individuals in the study can be used for the outcomes evaluation


\section{Weighting beyond the PS}
from \cite{Wang2019}:
Conventionally, the weights are estimated by modeling the propensities of receiving treatment or exhibiting missingness and then inverting the predicted propensities. However, with this approach it can be difficult to properly adjust for or balance the observed covariates. The reason is that this approach only balances covariates in expectation, by the law of large numbers, but in any particular data set it can be difficult to balance covariates, especially if the data set is small or if the covariates are sparse (Zubizarreta et al., 2011). In addition, this approach can result in very unstable estimates when a few observations have very large weights (e.g., Kang and Schafer 2007). To address these problems, a number of methods have been proposed recently. Instead of explicitly modeling the propensities of treatment or missingness, these methods directly balance the covariates. Some of these methods also minimize a measure of dispersion of the weights.

Most of these weighting methods balance covariates exactly rather than approximately. This is a subtle but important difference because approximate balance can trade bias for variance whereas exact balance cannot. Also, exact balance may not admit a solution whereas approximate balance may do so. For a fixed sample size, approximate balance may balance more functions of the covariates than exact balance.


\chapter{Balancing Weights}
  \section{Introduction}
  %hello

  \section{Estimating the Population Mean of Potential Outcomes}
  %

We study the following problem:
%%%%%%%%%%%%%%%%%%
% PROBLEM PRIMAL %
%%%%%%%%%%%%%%%%%%
\begin{problem}
  \begin{gather*}
    \underset{w_1, \ldots, w_n \in \R}{\mathrm{minimize}}
    \qquad
    \sum_{i = 1}^{n} T_i f(w_i)
  \end{gather*}
subject to the constraints
\begin{gather*}
    w_i T_i \ge 0,
    \qquad 
    i = 1, \ldots, n,
\end{gather*}
\begin{gather*}
  \sum_{ i = 1 }^{n}
    w_i T_i
  =
  1
\end{gather*}
\begin{gather*}
    \left| 
      \sum_{i = 1}^{n} w_i T_i B_k(X_i)
      - 
      \frac{1}{n} \sum_{i = 1}^{n} B_k(X_i)
    \right|
    \le 
    \delta_k,
    \qquad
    k = 1, \ldots, K
\end{gather*}
\end{problem}
%%%%%%%%%%%%%%%%%%%%%%%%%%%%%%%%%%%%%%%%%%%%%%%%%%%%%%%%%%%%%%%%%

We aim to prove that the solution to 
Problem \eqref{primal_weighting_binary}
is asymptotical consistent with the propensity score, i.e.

\begin{theorem}
  \label{ch_theorem_1}
  Under some (non-optimal) Assumptions,
  there exist constants
  $
    c_{1}
    , 
    c_{2}
    >
    0
  $
  and decreasing sequences
  $
    (\varepsilon_n^1)
    ,
    (\varepsilon_n^2)
    \subset
    (0,1]
  $
  that converge to 0
  such that
  for all
  $
    \tau
    \in
    (0,1]
  $
  there exists a constant
  $
    c_{\tau}
    \in
    [0,\infty)
  $
  only depending on 
  $ \tau $
  such that for all 
  $ n \ge 1 $
  and
  $
    \tau
    \in
    (0,1]
  $
  it holds
  \begin{align}
    \label{ch_theorem_1_alternative}
    \begin{split}
   & \P
    \left( 
  \norm{
    w^*_i
    -
    \frac{1}{n \pi(X_i)}
  }_\infty
  \le 
  c_{1}
  c_{\tau}
  \varepsilon_{n}^1 
   \right)
   \ge
  1 - \tau
  ,
\\
   &\quad \
   \norm{
    w^*_i
    -
    \frac{1}{n \pi(X_i)}
  }_{\P, 2}
  \le 
  c_{2}
  \varepsilon_n^{2},
      \end{split}
  \end{align}
  where 
  $w^*$
  is the solution to 
  Problem \eqref{primal_weighting_binary}.
\end{theorem}



%%%%%%%%%%%%%%%%%%%ASSUMPTION 1%%%%%%%%%%%%%%%%%%%%%%%%%%%%%%%%%%%%%%%


\begin{assumption}
  Assume, the following conditions hold:
  %
  \begin{subassumption}
    \label{assumption_1_i} 
    The minimizer 
    $
    \lambda_0 
    =
    \arg \min_{\lambda \in \Theta}
    \E
    \left[ 
      -T n 
      \rho 
      \left( 
      B(X)^T \lambda
      \right)
      +
      B(X)^T \lambda
    \right]
    $
    is unique,
    where 
    $\Theta \subseteq \R^n$ is the parameter space for $\lambda$.
  \end{subassumption}
  %
  \begin{subassumption}
    \label{assumption_1_ii} 
    The parameter space 
    $\Theta \subseteq \R^n$
    is compact.
  \end{subassumption}
  %
  \begin{subassumption}
    \label{assumption_1_iii}
    $\lambda_0 \in \text{int}(\Theta)$,
    where
    $\text{int}(\cdot)$
    stands for the interior of a set.
  \end{subassumption}
  %
  \begin{subassumption}
    \label{assumption_1_iv}
    There exists
    $\lambda^*_1 \in \Theta$
    such that
    $
      \norm{
        m^*(\cdot)
        -
        B(\cdot)^T \lambda^*_1
      }_\infty
      \le 
      \varphi_{m^*}
    $,
    where
    $
      m^*(\cdot)
      :=
      \left( \rho^{'} \right)^{-1}
      \left( \frac{1}{n \pi(\cdot)} \right).
    $
  \end{subassumption}
  %
  \begin{subassumption}
    \label{assumption_1_v}
    There exists a constant 
    $
      \varphi_{\pi} 
      \in 
      \left(0, \frac{1}{2} \right)
    $
    such that
    $
      \pi(x)
      \in 
      (
      \varphi_{ \pi },
      1 - \varphi_{ \pi }
      )
    $
    for all
    $
      x
      \in
      \mathcal{X}
    $
  \end{subassumption}
  %
  \begin{subassumption}
    \label{assumption_1_vi}
    There exists 
    $ \varphi_{\rho^{''}} > 0 $
    such that
    $ -\rho^{''} \ge \varphi_{\rho^{''}} > 0 $
  \end{subassumption}
  %
  \begin{subassumption}
    \label{assumption_1_vii}
    There exists 
    $ \varphi_{B(x) B(x)^T} > 0 $ 
    such that
    $
      B(x) B(x)^T 
      \succcurlyeq 
      \varphi_{B(x) B(x)^T} I 
    $
  \end{subassumption}
  %
  \begin{subassumption}
    \label{assumption_1_viii}
    There exists
    $ \varphi_{\norm{B}} > 0 $
    such that
    $
      \sup_{x \in \mathcal{X}} \norm{B(x)}_2
      \le 
      \varphi_{\norm{B}}
    $.
  \end{subassumption}
  %
  \begin{subassumption}
    \label{assumption_1_ix}
    The number of basis functions
    satisfies
    $
      K = o(n)
    $.
  \end{subassumption}
\end{assumption}


\section*{Plan of Proof}

It is easier to study the dual of 
Problem \eqref{primal_weighting_binary}.
Thus we employ results from convex analysis
\cite{Mordukhovich2022}
to establish
\begin{proposition}
  \label{ch_1_dual}
  The dual of Problem \eqref{primal_weighting_binary}
  is equivalent to the unconstrained optimization problem
  \begin{gather}
    \label{dual_weighting_binary}
    \underset{\lambda \in \R^K}{\text{minimize}}
    \quad
    \frac{1}{n}
    \sum_{j = 1}^{n} 
      \left[ 
        -T_j n 
        \rho 
        \left( 
          B(X_j)^T \lambda
        \right)
        +
        B(X_j)^T \lambda
      \right]
    +
    |\lambda|^T \delta,
  \end{gather}
  where 
  $
    B(X_j)=(B_k(X_j))_{1 \le k \le K}
  $
  denotes the $K$ basis functions of the covariates,
  $
    \rho(t) 
    :=
    \frac{t}{n}
    -
    t
    (h^{'})^{-1}(t)
    +
    h((h^{'})^{-1}(t))
  $
  with
  $
    h(x)
    :=
    f%
    \left( 
      \frac{1}{n}
      -
      x
    \right)
  $
  and
  $
    | \lambda | 
    :=
    \left(
    | \lambda_k | 
    \right)_{1 \le k \le K}
  $.
  Moreover,
  the primal solution
  $
    w^*_j
  $
  satisfies
  \begin{gather}
    \label{ch_1_w*}
    w^*_j
    =
    \rho^{'}
    \left(
      B(X_j)^T
      \lambda^\dagger
    \right)
  \end{gather}
  for 
  $
    j
    =
    1, \ldots , n
  $,
  where 
  $
    \lambda^\dagger
  $
  is the solution to the dual optimization problem.
\end{proposition}
%%%%%%%%%%%%%%%%%%%%%%%%%%%%%%%%%%%%%%%%%%%%%%%%%%%%%%%%%%%%%%%%%%%%%
%%%%%%%%%%%%%%%%%%%%%%%%%%%%%%%%%%%%%%%%%%%%%%%%%%%%%%%%%%%%%%%%%%

The core of the subsequent analysis is based on
Assumption~\ref{assumption_1_iv}, 
i.e.
the existence of an oracle parameter
$\lambda_1^{*}$
in a sieve estimate of the true propensity score (or a transformation). 
It is then natural to enquire about the convergence of the dual solution
$
  \lambda^\dagger
$
to 
$
  \lambda^*_1
$.
Making certain assumptions and employing
matrix concentration inequalitys \cite{Tropp2015}
we can establish
\begin{proposition}
  \label{ch_1_near_oracle}
  Under some (non-optimal) Assumptions,
  there exists a constant
  $
    c_{3}
    >
    0
  $
  and a decreasing sequence
  $
    (\varepsilon_n^3)
    \subset
    (0,1]
  $
  that converges to 0
  such that
  for all
  $
    \tau
    \in
    (0,1]
  $
  there exists a constant
  $
    \tilde{
    c_{\tau}
    }
    \in
    [0,\infty)
  $
  only depending on 
  $ \tau $
  such that for all 
  $ n \ge 1 $
  and
  $
    \tau
    \in
    (0,1]
  $
  it holds
  \begin{gather}
    \P
    \left( 
      \norm{
        \lambda^\dagger
        -
        \lambda^*_1
      }_2
      \le
      c^3
      \tilde{
        c_{\tau}
      }
      (
        \varepsilon_n^3
      )
    \right)
    \ge 
    1 - \tau.
  \end{gather}
\end{proposition}%
It is then straightforward to prove a more general result
than
Theorem~\ref{ch_theorem_1}%
.
%
\begin{theorem}
  Under some (non-optimal) Assumptions,
  there exist constants
  $
    c_{1}
    , 
    c_{2}
    >
    0
  $
  and decreasing sequences
  $
    (\varepsilon_n^1)
    ,
    (\varepsilon_n^2)
    \subset
    (0,1]
  $
  that converge to 0
  such that
  for all
  $
    \tau
    \in
    (0,1]
  $
  there exists a constant
  $
    c_{\tau}
    \in
    [0,\infty)
  $
  only depending on 
  $ \tau $
  such that for all 
  $ n \ge 1 $
  and
  $
    \tau
    \in
    (0,1]
  $
  it holds
  \begin{align*}
   & \P
    \left( 
  \norm{
    w^*(\cdot)
    -
    \frac{1}{n \pi(\cdot)}
  }_\infty
  \le 
  c_{1}
  c_{\tau}
  \varepsilon_{n}^1 
   \right)
   \ge
  1 - \tau
  ,
\\
   &\quad \
   \norm{
    w^*(X)
    -
    \frac{1}{n \pi(X)}
  }_{\P, 2}
  \le 
  c_{2}
  \varepsilon_n^{2},
  \end{align*}
  where 
  $w^*(X)$
  is as in \eqref{ch_1_w*}
  without the index.
\end{theorem}

\section*{Proof of theorem 2.2}


\begin{proof}
  Motivated by Proposition~\ref{syu_1_result}
  we consider
 \begin{gather}
   \label{ch_1_G}
   G(\lambda)
   :=
      \frac{1}{n}
      \sum_{j = 1}^{n} 
      \left[ 
        -T_j n 
        \rho 
        \left( 
          B(X_j)^T \lambda
        \right)
      +
      B(X_j)^T \lambda
      \right]
      +
      |\lambda|^T \delta.
 \end{gather} 
 Since 
 $\rho \in C^2(\R)$
 we can employ 
 \eqref{ch_1_G},
 Corollary~\ref{syu_taylor_corollary}
 and
 Proposition~\ref{syu_triangle}
 to get
 \begin{align}
   \label{ch_1_proof1_1}
  \begin{split}
  & 
  G(\lambda^*_1 + \Delta) 
  -
  G(\lambda^*_1)
  \\
  &
  \quad
  \ge
      \frac{1}{n}
      \sum_{j = 1}^{n} 
      \left[ 
        -T_j n 
        \rho^{'} 
        \left( 
          B(X_j)^T \lambda^*_1
        \right)
      +
      1
      \right]
      \Delta^T B(X_j)
      \\
  & 
  \qquad
  +
      \frac{1}{2}
      \sum_{j = 1}^{n} 
        -T_j  
        \rho^{''} 
        \left( 
          B(X_j)^T (\lambda^*_1 + \xi \Delta)
        \right)
        \Delta^T
        \left( 
          B(X_j)
          B(X_j)^T
        \right)
        \Delta
        \\
  &
  \qquad
  -
      |\Delta|^T \delta
  \\
  &  
  \quad
  \ge
    - \norm{\Delta}_2
    \left( 
    \norm{
      \frac{1}{n}
      \sum_{j = 1}^{n} 
      \left[ 
        -T_j n 
        \rho^{'} 
        \left( 
          B(X_j)^T \lambda^*_1
        \right)
      +
      1
      \right]
      B(X_j)
    }_2
    +
    \norm{\delta}_2
    \right)
    \\
  & 
  \qquad
  +
  n
  \norm{\Delta}^2_2
   \varphi_{\rho^{''}}
  \underline{\varphi_{BB^T}}
  \\
  &
  \quad
  :=
  -\norm{\Delta}_2
  (I_1 + \norm{\delta}_2)
  +
  \norm{\Delta}^2_2
  I_2
  .
  \end{split}
\end{align}
The second inequality is due to 
the Cauchy-Schwarz-Inequality 
and
Assumptions~\ref{assumption_1_vi} and \ref{assumption_1_vii}
.
We want to establish probabilistic upper bounds of the factor associated with 
$
  -\norm{\Delta}_2
$. 
This will be done with appropriate assumptions on
$
  \norm{\delta}_2
$
and a thorough analysis of $I_1$.
If we then restrict lower bounds of
$
  I_2
$
to approprietely slow convergence to 0,
e.g. by assumptions on
$
\varphi_{\rho^{''}}
$
and
$
\underline{
  \varphi_{BB^T}
}
$,
we can choose
$
  \norm{\Delta}_2
$
large enough,
such that 
\eqref{ch_1_proof1_1}
yields
$
  G(\lambda^*_1 + \Delta) 
  -
  G(\lambda^*_1)
  >
  0
$
with arbitrarily large probabitity for $n$ large enough.
With Proposition~\ref{syu_1_result} it follows then immediately
Proposition~\ref{ch_1_near_oracle}
.
\subsection*{Analysis of $I_1$}
We want to use Assumption~\ref{assumption_1_iii}.
Thus we perform the following split:
\begin{align}
  \label{ch_1_i1}
  \begin{split}
  I_1 
  &\le
    \norm{
      \sum_{j = 1}^{n} 
        T_j  
      \left[ 
        \rho^{'} 
        \left( 
          B(X_j)^T \lambda^*_1
        \right)
      -
      \frac{1}{n \pi(X_j)}
      \right]
      B(X_j)
    }_2
  \\
  &+
    \norm{
      \frac{1}{n}
      \sum_{j = 1}^{n} 
      \left[ 
        \frac{T_j}{\pi(X_j)}
      -
      1
      \right]
      B(X_j)
    }_2
    \\
  &=:
  J_1 + J_2
  \end{split}
\end{align}

\subsubsection*{Analysis of $J_1$}

By the Lipschitz-continuity of 
$\rho^{'}$,
Assumption~\ref{assumption_1_viii}
and
Assumption~\ref{assumption_1_iv},
$T \in \{0, 1\}$
and 
the triangle inequality 
we have
\begin{gather}
  \label{ch_1_j1_result}
  J_1 
  \le
  n L_{\rho^{'}}\varphi_{\norm{B(x)}} \varphi_{m^*}
\end{gather}

\subsubsection*{Analysis of $J_2$}
We want to employ 
Theorem~\ref{rmineq_bernstein}%
. %
To this end we define
the independent random matrices
\begin{align}
  \label{ch_1_j2_1}
  \begin{split}
  A_j
  &:=
      \frac{1}{n}
      \left[ 
        \frac{T_j}{\pi(X_j)}
      -
      1
      \right]
      B(X_j)
      \,,
      \quad 
      j=1,\ldots,n,
\\
  S
  &:=
  \sum_{j = 1}^{n} 
  A_j
  \end{split}
\end{align}
and check conditions 
\eqref{rmineq_bernstein_cond_1} 
and
\eqref{rmineq_bernstein_cond_2}.
Note that 
$
  \norm{S}_2
  =
  J_2
$
.
By the properties of conditional expectation 
it holds
\begin{gather}
  \label{ch_1_j2_2}
  \E
  \left[  
    \frac{T_j}{\pi(X_j)}
    B(X_j)
  \right]
  =
  \E
  \left[  
    \E
    \left[  
      T_j
      \, | \,
      X_j
    \right]
    \frac{1}{\pi(X_j)}
    B(X_j)
  \right]
  =
  \E[B(X_j)]. 
\end{gather}
Taking the expectation in 
\eqref{ch_1_j2_1} 
and using
\eqref{ch_1_j2_2}
we get
$\E[A_j] = 0$
for all 
$
  j=1,\ldots,n
$
.
Since
\begin{gather}
  \label{ch_1_j2_3}
  \left| 
      \frac{T_j}{\pi(X_j)}
      -
      1
      \right|
  \le
  1 + \frac{1 - \varphi_{\pi}}{\varphi_{\pi}}
  =
  \frac{1}{\varphi_{\pi}}
\end{gather}
by Assumption~\ref{assumption_1_v},
we can employ Assumption~\ref{assumption_1_viii}
together with 
\eqref{ch_1_j2_3}
and
\eqref{ch_1_j2_1}
to get
\begin{gather}
  \label{ch_1_j2_4}
  \norm{A_j}_2
  \le
  \frac{\varphi_{\norm{B}}}{n \varphi_{\pi}}
  =:
  L.
\end{gather}
Thus,
condition \eqref{rmineq_bernstein_cond_1}
is satisfied.
Next we turn to the matrix variance statistic
$ v(S) $
\eqref{rmineq_bernstein_cond_2}.
By 
\eqref{ch_1_j2_1}
and
\eqref{ch_1_j2_3}
we have
\begin{gather}
  \label{ch_1_j2_5}
  \E
  \left[ 
    A_j A_j^T
  \right]
  \le
  \left( 
    \frac{1}{n \varphi_{\pi}}
  \right)^2
  \E
  \left[ 
    B(X) B(X)^T
  \right]
\end{gather}
and by \eqref{ch_1_j2_4} 
\begin{gather}
  \label{ch_1_j2_6}
  \E
  \left[ 
    A_j^T A_j
  \right]
  \le
  L^2.
\end{gather}
Since 
$
  \max\{a,b\}
  \le
  | a | + | b |
$
we can use 
\eqref{ch_1_j2_5}
and
\eqref{ch_1_j2_6}
to get
\begin{gather}
  \label{ch_1_j2_7}
  v(S)
  \le
  \frac{1}{n}
  \frac{\lambda_{\max}}{\varphi_{\pi}^2}
  +
  n L^2,
\end{gather}
where 
$\lambda_{\max}$ 
is the maximal eigenvalue of
the symmetric (non-random) matrix
$
  \E
  \left[ 
    B(X) B(X)^T
  \right]
$.
Having dealt with 
\eqref{rmineq_bernstein_cond_1} 
and
\eqref{rmineq_bernstein_cond_2}
we can establish 
the
expectation bound
\eqref{rmineq_bernstein_expectation_bound}
of Theorem~\ref{rmineq_bernstein}.
Together with 
\eqref{ch_1_j2_4}
and
\eqref{ch_1_j2_7}
we get
\begin{align}
  \label{ch_1_j2_8}
  \begin{split}
    &\E[J_2]
\\
  &\le
  \sqrt{
    \frac{
    2 \log (K + 1)
    \left( 
      \lambda_{\max}  + \varphi_{\norm{B}}^2
    \right)
    }
    {
      n \varphi_{\pi}^2
    }
  }
  +
  \frac{
    \log (K + 1)
    \varphi_{\norm{B}}
  }
  {
    3 n \varphi_{\pi}
  }
\\
  &\le
  \frac{1}{\varphi_{\pi}}
  \sqrt{
    \frac{
      \log(K + 1)
    }
    {
      n
    }
  }
  \left[
    \varphi_{\norm{B}}
    \left(
      \sqrt{2}
      +
      \frac{1}{3}
      \sqrt{
        \frac{\log (K + 1)}{n}
      }
    \right)
    +
    \sqrt{
      2
      \lambda_{\max}
    }
  \right].
\end{split}
\end{align}
Since
$
  K = o(n)
$
by Assumption~\ref{assumption_1_ix}
we can discuss the other influences on the quality of the bound
\eqref{ch_1_j2_8}.
On a high-level it is readily clear
that appropriate bounds on
$
\varphi_{\pi},
\varphi_{\norm{B}}
$
and
$
  \lambda_{\max}
$
will shrink
$
  \E[J_2]
$
to 0 
and will assist in establishing learning rates.

We could also have invoked the probabitity bound 
\eqref{rmineq_bernstein_probability_bound} 
of Theorem~\ref{rmineq_bernstein}.
But for the sake of simplicity we prefer the combination 
of the expectation bound 
\eqref{ch_1_j2_8}
and the Markov inequality.
With the latter we get
\begin{align}
  \label{ch_1_j2_result}
  J_2
  \le
  \frac{1}{\tau}
  \frac{1}{\varphi_{\pi}}
  \sqrt{
    \frac{
      \log(K + 1)
    }
    {
      n
    }
  }
  \left[
    \varphi_{\norm{B}}
    \left(
      \sqrt{2}
      +
      \frac{1}{3}
      \sqrt{
        \frac{\log (K + 1)}{n}
      }
    \right)
    +
    \sqrt{
      2
      \lambda_{\max}
    }
  \right]
\end{align}
with probabitity 
$
  \ge
  1 - \tau
$.

%%%%%%%%%%%%%%%%%%%%%%%%%%%%%%%%%%%%%


If we choose
$
\norm{\Delta}_2
$
to be
\begin{align}
  \label{ch_1_choice_norm_Delta}
  \begin{split} 
    &\left(
     \sqrt{2}
    \frac{1}{\tau}
    \frac{1}{\varphi_{\pi}}
    \sqrt{
      \frac{
        \log(K + 1)
      }
      {
        n^3
      }
    }
    \left[
      \varphi_{\norm{B}}
      \left(
        1
        +
        \sqrt{
          \frac{\log (K + 1)}{n}
        }
      \right)
      +
      \sqrt{
        \lambda_{\max}
      }
  \right]
  \right.
\\
  &
  \quad
  +
     L_{\rho^{'}}
     \varphi_{\norm{B}}
     \varphi_{m^*}
  +
  \left.
  \frac{
    \norm{\delta}_2
  }
  {n}
  \right)
  \frac{1}
  {
    \varphi_{\rho^{''}}
    \underline{
      \varphi_{BB^T}
    }
  }
\end{split}
\end{align}
we get by
\eqref{ch_1_proof1_1},
\eqref{ch_1_i1},
\eqref{ch_1_j1_result},
\eqref{ch_1_j2_result}
and 
Proposition~\ref{syu_1_result}
\begin{align}
  \label{ch_1_near_oracle_probabilistic_bound}
  \begin{split}
  \P
  \left( 
      \norm{
        \lambda^\dagger
        -
        \lambda^*_1
      }_2
      \le
      C 
  \right)
  &=
  \P
  \left( 
    \inf_{\norm{\Delta}_2 = C}
    G(\lambda^*_1 + \Delta) 
    -
    G(\lambda^*_1)
    >
    0
  \right)
  \\
  &\ge
  1 - \tau,
  \end{split}
\end{align}
where 
$
  C
$
is as in 
\eqref{ch_1_choice_norm_Delta}.
With appropriate Assumptions (as discussed before) we can then establish
Proposition~\ref{ch_1_near_oracle}.

We can invoke 
\eqref{ch_1_near_oracle_probabilistic_bound}
to derive bounds as in 
Theorem~\ref{ch_theorem_1_alternative}:
\begin{align*}
  \norm{
    w^*(X)
    -
    \frac{1}{n \pi(X)}
  }_{\P, 2}
  &\le
  L_{\rho^{'}}
  \left[ 
    \norm{
      B(X)^T 
      \left(
        \lambda^\dagger
        -
        \lambda^*_1
      \right)
    }_{\P,2}
    \right.
    \\
    &+
    \left. 
    \norm{
      m^*(X)
      -
      B(X)^T 
      \lambda^*_1
    }_{\P,2}
  \right]
  \\
  &\le
  L_{\rho^{'}}
  \left(
    \varphi_{\norm{B}}
    \sqrt{
      C^2(1 - \tau)
      +
      \text{diam}(\Theta)^2
      \tau
    } 
    +
    \varphi_{m^*}
  \right)
\end{align*}
%%%%%%%%%%%%%%%%%%%
\begin{align*}
  \norm{
    w^*(\cdot)
    -
    \frac{1}{n \pi(\cdot)}
  }_{\infty}
  &\le
  L_{\rho^{'}}
  \left[ 
    \norm{
      B(\cdot)^T 
      \left(
        \lambda^\dagger
        -
        \lambda^*_1
      \right)
    }_{\infty}
    \right.
    \\
    &+
    \left. 
    \norm{
      m^*(\cdot)
      -
      B(\cdot)^T 
      \lambda^*_1
    }_{\infty}
  \right]
  \\
  &\le
  L_{\rho^{'}}
  \left(
    \varphi_{\norm{B}}
    C
    +
    \varphi_{m^*}
  \right)
\end{align*}
with probabitity greater than 
$1 - \tau$.
\end{proof}
\begin{remark}
  By Corollary~\ref{rmineq_intrinsic_bernstein_vector_corollary}
  we can get rid of the $\log(K)$ term in
\eqref{ch_1_choice_norm_Delta}.
\end{remark}
\begin{remark}
  By the matrix Rosenthal-Pinelis Inequality
  \cite{Chen2012}[Thm.A.1]
  we can weaken Assumption~\ref{assumption_1_v} to a lower bound on the expectition of 
  $\pi(X)$
\end{remark}
The next step consists of strenghtening the Assumptions to get concrete learning rates.
This can be done in a series of examples.

  \section{Application of Convex Optimization}
  \begin{assumption}
  \begin{enumerate}[label={(\roman*)}]
    Assume that the map 
    $
      f: \R \to \overline{\R}
    $
    has the following properties.
    \item
      $
        f 
        \ 
        \text{is strictly convex.}
      $
    \item
      $
        f
        \ 
        \text{is lower-semicontinuous and continuously differentiable on}
        \ 
        \mathrm{int}(\mathrm{dom}(f))
        .
      $
    \item
      $
        \text{The derivative of}\ 
        f\ 
        \text{on}\ 
        \mathrm{int}(\mathrm{dom}(f))
        \ 
        \text{is a diffeomorphism.}
      $
    \item
      $
        \text{The Legendre transformation}
        \ 
        f^*
        \ 
        \text{of}\ 
        f
        \ 
        \text{is finite}
        .
      $
    \item
      $
        \text{The function}\
        x\mapsto xt - f(x)
        \ 
        \text{takes its supremum on}
        \ 
        \mathrm{int}(\mathrm{dom}(f))
        \ 
        \text{for all}
        \ 
        t\in \R.
      $
  \end{enumerate}
\end{assumption}

We consider the following optimization problem.
\begin{problem}
  \begin{gather*}
    \underset{w_1, \ldots, w_n \in \R}{\mathrm{minimize}}
    \qquad
    \sum_{i = 1}^{n} T_i f(w_i)
  \end{gather*}
subject to the constraints
\begin{gather*}
    w_i T_i \ge 0,
    \qquad 
    i = 1, \ldots, n,
    \\
  \sum_{ i = 1 }^{n}
    w_i T_i
  =
  1
  \\
    \left| 
      \sum_{i = 1}^{n} w_i T_i B_k(X_i)
      - 
      \frac{1}{n} \sum_{i = 1}^{n} B_k(X_i)
    \right|
    \le 
    \delta_k,
    \qquad
    k = 1, \ldots, K
\end{gather*}
\end{problem}
\pagebreak
~\begin{ftheorem}
  Under Assumption,
  the dual of the above Problem is the unconstrained optimization problem 
  \begin{gather*}
    \underset{\lambda \in \R^K}{\mathrm{minimize}}
    \qquad
    \frac{1}{n}
    \sum_{i = 1}^{n} 
    nT_i 
    f^*
    (
      \inner{B(X_i)}{\lambda}
    )
    -
    \inner{B(X_i)}{\lambda}
    +
    \inner{\delta}{\left| \lambda \right|}
    ,
  \end{gather*}
  where
  $
  t \mapsto
    f^*(t)
    =
    t\,(f^{'})^{-1}(t)
    -
    f
    \left( 
      (f^{'})^{-1}(t)
    \right)
  $
  is the Legendre transformation of $f$,
  $
    B(X_i)
    =
    \left[ 
      B_1(X_i)
      ,
      \ldots
      ,
      B_K(X_i)
    \right]
    ^\top
  $
  denotes the $K$ basis functions of the covariates 
  of unit $i\in \left\{ 1, \ldots, n \right\}$
  and
  $
    \left| \lambda \right|
    =
    \left[ 
      \left| \lambda_1 \right|
      ,
      \ldots
      ,
      \left| \lambda_K \right|
    \right]
    ^\top
    ,
  $
  where $\left| \,\cdot\, \right|$
  is the absolute value of a real-valued scalar.
  Moreover, if $\lambda^\dagger$
  is an optimal solution then
  \begin{gather}
    w_i^*
    =
    (f^{'})^{-1}
    (
      \inner{B(X_i)}{\lambda^\dagger}
    )
    ,
    \quad
    i\in \left\{ 1,\ldots, n \right\}
  \end{gather}
  are the unique optimal solutions to (P)
  .
\end{ftheorem}
\begin{proof}
  We prove the following Lemma at the end of the section.
  \begin{lemma}
    The dual of the optimization problem is
  \begin{gather*}
    \underset{\lambda \in \R^{2K}}{\mathrm{minimize}}
    \qquad
    \frac{1}{n}
    \sum_{i = 1}^{n} 
    nT_i 
    f^*
    (
    \inner{Q_{\bullet i}}{\lambda}
    )
    -
    \inner{Q_{\bullet i}}{\lambda}
    +
    \inner{d}{\lambda}
  \end{gather*}
  subject to
  \begin{gather}
    \lambda_k \ge 0
    \quad
    \text{for all}\ 
    k \in \left\{ 1, \ldots, K \right\}
    ,
  \end{gather}
  where
  \begin{gather}
    \mathbf{Q}
    :=
    \begin{bmatrix}
      \mathbf{I}_n\\
      \mathbf{B}(\mathbf{X})\\
      -\mathbf{B}(\mathbf{X})
    \end{bmatrix}
    ,
    \qquad
    \mathbf{B}(\mathbf{X})
    :=
    \begin{bmatrix}
      B(X_1), \ldots, B(X_n)
    \end{bmatrix}
    ,
    \qquad
    \text{and}
    \qquad
    d
    :=
    \begin{bmatrix}
      0_n\\
      \delta \\
      \delta
    \end{bmatrix}
    .
  \end{gather}

  \end{lemma}

\end{proof}

  \section{Application of Matrix Concentration Inequalities}
  %\subsubsection*{
  Analysis of 
  $
  \E[\max_{i \le r}\norm{\mathbf{A}_i}^2]
  $
}
We have
\begin{gather}
  \mathbf{A}_i
  :=
  \frac{1}{r}
  \left( 
    \frac{1-\pi_i}{\pi_i}
  \right)
  \mathbf{B}(X_i)
  \qquad
  \text{for}
  \ 
  i\in \{1, \ldots, r\}.
\end{gather}
Since we take the maximum over a finite set it is attained for some 
$i^*\in \left\{ 1, \ldots, r \right\}$:
\begin{align}
  \begin{split}
  \E[\max_{i \le r}\norm{\mathbf{A}_i}^2]
  &=
  \E[\norm{\mathbf{A}_{i^*}}^2]
  \\
  &=
  \frac{1}{n^2}
  \E \left[ 
    \left(
      \frac{1-\pi_{i^*}}{\pi_{i^*}}
    \right)^2
    \norm{\mathbf{B}(X_{i^*})}^2
  \right]
  \le
  \frac{1}{n^2}
  \E \left[ 
    \left(
      \frac{1-\pi_{i^*}}{\pi_{i^*}}
    \right)^4
  \right]^\frac{1}{2}
  \E[
    \norm{\mathbf{B}(X_{i^*})}^4
  ]^\frac{1}{2}
  \\
  &\le
  \frac{K}{n^2}\sqrt{C_\pi C_\mathbf{B}}
\end{split}
\end{align}
In the last two steps we applied the Cauchy-Schwarz inequality
and Assumption
\begin{assumption}
  There exists $C_\pi \ge 1$ such that
  $
    \E
    \left[ 
    \left(
      \frac{1-\pi_{i}}{\pi_{i}}
    \right)^4
    \right]
    \le C_\pi
  $
  for all $i\in \left\{ 1, \ldots, r \right\}$
  .
\end{assumption}
 
\begin{remark}
  If we assume a logistic regression model for the propensity score
  it holds for some $\theta \in \R^N$ ($N$ is the number of covariates)
  \begin{gather}
    \label{rmineq_rp_1}
    \frac{1-\pi(X)}{\pi(X)}
    =
    \exp(-\theta X)
    \qquad
    \text{and}
    \qquad
    \E
    \left[ 
    \left(
    \frac{1-\pi(X)}{\pi(X)}
    \right)^4
    \right]
    =
    \E
    [
    \exp(-4\theta X)
    ]
    =
    M_X(-4\theta)
    ,
  \end{gather}
  where $M_X$ is the momement-generating function of $X$.
   While the first quantity in \eqref{rmineq_rp_1}
   may be unbounded when $X$ has unbounded support, the latter quantity in \eqref{rmineq_rp_1} is still bounded for reasonable choices of $X$.
\end{remark}
\begin{assumption}
  There exists $C_\mathbf{B} \ge 1$ such that
  $
  \E[
    \mathbf{B}_k(X_i)^4
  ]
  \le C_\mathbf{B}
  $
  for all $(k,i)\in \left\{ 1, \ldots, K \right\}\times \left\{ 1, \ldots, r \right\}$
  .
\end{assumption}
\begin{remark}
With Assumption we also get a bound on the fourth moment of 
  $
  \norm{\mathbf{B}(X_{i})}
  $. Indeed, by the convexity of 
  $x\mapsto x^2$, the monotonicity and linearity of the expectation it holds   
  \begin{align}
    \begin{split}
  \E[
  \norm{\mathbf{B}(X_{i})}^4
  ] 
  &=
  \E
  \left[ 
    \left( 
      \sum_{k=1}^{K}
      \mathbf{B}_k^2(X_i)
    \right)^2
  \right]
  =
  K^2
  \E
  \left[ 
    \left( 
      \sum_{k=1}^{K}
      \frac{1}{K}
      \mathbf{B}_k^2(X_i)
    \right)^2
  \right]
  \le
  K^2
  \E
  \left[ 
      \sum_{k=1}^{K}
      \frac{1}{K}
      \mathbf{B}_k^4(X_i)
  \right]
  \\
  &=
  K
  \sum_{k=1}^{K}
  \E
  \left[ 
      \mathbf{B}_k^4(X_i)
  \right]
  \le
  K^2 C_\mathbf{B}
  \end{split}
  \end{align}
\end{remark}



\chapter{Convex Analysis}
In our application we want to analyze a convex optimization problem by its dual problem.
In particular we want to obtain primal optimal solutions from dual solutions.
To accomplish the task we need technical tools from convex analysis, 
mainly conjugate calculus and some KKT related results.

Our starting point is the support function intersection rule, which we will prove in full detail employing a theorem on convex separation is finite dimensions.
To this end, we will have a closer look in relative interiors and support functions.
As an application we may prove the conjugate chain and sum rule,
which are vital to application of duality. As a simple corollary we will obtain the classical Fenche-Rockafellar Duality theorem
which gives general conditions for dual und primal optimal values to be equal.

The material we present is very well known, so we claim no originality. We orient our exposition closly by
\cite{Rockafellar1970, Mordukhovich2022}
.

We finish the chapter with an exposition of \cite{Tseng1991}, where for strictly convex functions we get a dual relationship in terms of the optimal solutions.
  \section{A Convex Analysis Primer}
  Excursively, we present some well known definitions and facts from
convex analysis. For details, see, e.g., \cite{Mordukhovich2022}.

A subset $C\subseteq \R^n$ is called \textbf{convex set}, 
if for all $x,y\in C$ and all $\lambda\in [0,1]$,
we have 
$
  \lambda x + (1-\lambda)y 
  \in
  C
  .
$
The Cartesian product of convex sets is convex. The intersection of a collection of convex sets is also convex.




Given (not necessary convex) sets 
$
  \Omega,
  \Omega_1,
  \Omega_2
  \subseteq 
  \R^n
$
and
$\lambda\in\R,$
define the \textbf{set addition} and \textbf{multiplication}
by a real scalar as 
$
  \Omega_1 
  +
  \Omega_2 
  :=
  \left\{ 
    x_1 + x_2 
    \colon
    x_1 \in \Omega_1
    ,
    x_2 \in \Omega_2
  \right\}
$
and
$
  \lambda \Omega
  :=
  \left\{ 
    \lambda x
    \colon
    x\in\Omega
  \right\}.
$
For convex sets the addition and multiplication by a real scalar are convex.

Throughout this section, we shall denote by 
$
  B
  :=
  \left\{ 
    x=
    [x_1, \ldots, x_n]^\top
    \in\R^n
    \colon
    (
    \sum_{i=1}^{n} 
    x_i^2
    )
    ^{1/2}
    \le 1
  \right\}
$
\todo[color=green!40,inline]{Solve editorial issue with ball.}
the \textbf{Euclidian unit ball} in $\R^n.$
This is a closed convex set. For any $a\in\R^n,$ 
the \textbf{ball with radius $\varepsilon >0$ and center $a$}
is given by
$
  \left\{ 
   a+x 
    \in\R^n
    \colon
    (
    \sum_{i=1}^{n} 
    x_i^2
    )
    ^{1/2}
    \le \varepsilon
  \right\}
  =
  a
  +
  \varepsilon B
  .
$
For any set $\Omega$ in $\R^n,$ 
the set of points $x$ whose distance from $\Omega$
does not exceed $\varepsilon$ is 
$
\Omega + \varepsilon B.
$
The \textbf{closure} $\mathrm{cl}(\Omega)$
and \textbf{interior} $\mathrm{int}(\Omega)$
of $\Omega$
can therefore be expressed by 
$
  \mathrm{cl}(\Omega)
  =
  \bigcap_{\varepsilon>0}
  \Omega + \varepsilon B
$
and
$
  \mathrm{int}(\Omega)
  =
  \left\{ 
    x\in \Omega
    \colon
    \text{there exists $\varepsilon>0$ such that }
    x+\varepsilon B
    \subseteq
    \Omega
  \right\}
  .
$

  A set 
  $A\subseteq \R^n$
  is called \textbf{affine set}, if
  $
    \alpha x + (1-\alpha)y \in A
    \quad
    \text{for all}
    \ 
    x,y \in A
    \ 
    \text{and}
    \ 
    \alpha \in \R.
  $
  The \textbf{affine hull} 
  $\mathrm{aff}(\Omega)$
  of a set 
  $\Omega\subseteq \R^n$
  is the smallest affine set that includes $\Omega.$
A mapping 
$
  A: \R^n \to \R^m
$
is called \textbf{affine mapping} if there exist a linear mapping
$
  L: \R^n \to \R^m
$
and a vector $b\in \R^m$
such that
$
  A(x)
  =
  L(x)
  +
  b
  \ 
  \text{for all}\ 
  x \in \R^n
  .
$
The image and inverse image/preimage of convex sets under affine mappings are also convex.





Because the notion of interior is not precise enough for our purposes
we define the relative interior which is the interior relative to the affine hull.
This concept is motivated by the fact that a line segment embedded in $\R^2$
does have a natural interior in $\R$ which is not a true interior in $\R^2.$
The relative interior of $C$ is defined as the interior which results when
$C$ is regarded as a subset of its affine hull.
\begin{definition}
  Let 
  $\Omega\subseteq \R^n.$
  We define the \textbf{relative interior} of $\Omega$ by
  \begin{gather}
    \mathrm{ri}(\Omega)
    :=
    \left\{ 
      x \in \Omega 
      \colon
      \text{there exists}\ 
      \varepsilon > 0\ 
      \text{such that}\ 
      (
        x+\varepsilon B
      )
      \cap
      \mathrm{aff}(\Omega)
      \subset
      \Omega
    \right\}.
  \end{gather}
\end{definition}
Next we collect some useful properties of relative interiors.
\begin{proposition}
  Let $C$ be a non-empty convex set in $\R^n.$ Then we get the representation
\begin{enumerate}[label={(\roman*)}]
  \item
    $
    \mathrm{ri}(C)
    =
    \left\{ 
      z \in C
      \colon
      \text{for all}\ 
      x \in C \ 
      \text{there exists}\ 
      t > 0 \ 
      \text{such that}\ 
      z + t (z-x)
      \in C
    \right\}.
    $
  \item
    $
      \mathrm{ri}(C)
      \neq
      \emptyset
      \ 
      \text{if}\ 
      C\neq \emptyset
      .
    $
  \item
    $
      \mathrm{cl}(C)
      \ 
      \text{and}\ 
      \mathrm{ri}(C)
    $
    are convex sets.
  \item
    $
      \mathrm{cl}(\mathrm{ri}(C))
      =
      \mathrm{cl}(C)
      \ 
      \text{and}
      \  
      \mathrm{ri}(\mathrm{cl}(C))
      =
      \mathrm{ri}(C)
      .
    $
  \item
    Suppose
    $
      \bigcap_{i\in I} C_i
      \neq
      \emptyset
    $
    for a finite index set $I.$
    Then
    $
      \mathrm{ri}
      \left( 
        \bigcap_{i\in I} C_i
      \right)
      =
      \bigcap_{i\in I}  
      \mathrm{ri}(C_i)
      .
    $
    \item
      Let 
      $
        L:\R^n \to \R^m
      $
      be a linear mapping. Then
      $
        \mathrm{ri}(L(C))
        =
        L(\mathrm{ri}(C))
        .
      $
      If additionally it holds
      $
        L^{-1}(\mathrm{ri}(C))
        \neq
        \emptyset
      $
      we have
      $
      \mathrm{ri}(L^{-1}(C))
        =
        L^{-1}(\mathrm{ri}(C))
        .
      $
      \item
        $
          \mathrm{ri}(C_1\times C_2)
          = 
          \mathrm{ri}(C_1)
          \times
          \mathrm{ri}(C_2).
        $
      \item
        $
          \mathrm{ri}(C_1)
          \cap
          \mathrm{ri}(C_2)
          =
          \emptyset
        $
        if and only if
        $
          0 \notin
          \mathrm{ri}
          (C_1 - C_2)
          .
        $
\end{enumerate}

\end{proposition}

\todo[color=green!40,inline]{Order results to give pretty proof.}

\begin{proof}
\begin{enumerate}[label={(\roman*)}]
  \item
    \cite[Theorem~6.4]{Rockafellar1970}
  \item
    \cite[Theorem~6.2]{Rockafellar1970}
  \item
    \cite[Theorem~6.2]{Rockafellar1970}
  \item
    \cite[Theorem~6.3]{Rockafellar1970}
  \item
    \cite[Theorem~6.5]{Rockafellar1970}
  \item
    \cite[Theorem~6.6-6.7]{Rockafellar1970}
  \item
Let
  $
  (z_1, z_2)
  \in 
  \mathrm{ri}(C_1\times C_2).
  $
  Then for all 
  $
  (x_1, x_2)
  \in 
  C_1\times C_2
  $
  there exists
  $t>0$
  such that
  \begin{gather}
      z_i + t (z_i-x_i)
      \in C_i
      \qquad
      \text{for}\ 
      i\in \left\{ 1,2 \right\}.
  \end{gather}
  This proves $\subseteq.$
  Suppose 
  $
    z_1 
    \in
    \mathrm{ri}(C_1)
  $
  and
  $
    z_2 
    \in
    \mathrm{ri}(C_2).
  $
  Let
  $
    (x_1,x_2)\in C_1\times C_2
  $
  with
  If $t_1=t_2$ everything is clear.
  W.l.o.g.
  assume 
  $t_1<t_2$. Define $\theta:=\frac{t_1}{t_2}\in (0,1).$
  By the convexity of $C_2$ it follows
  \begin{gather}
    z_2 + t_1 (z_2 - x_2)
    =
    \theta
    (
    z_2 + t_2 (z_2 - x_2)
    )
    +
    (1-\theta)z_2
    \in C_2.
  \end{gather}
  Thus
  $
  (z_1,z_2)\in\mathrm{ri}(C_1\times C_2). 
  $
  This proves $\supseteq$ and equality.
  \item
    \cite[Theorem~2.92]{Mordukhovich2022}
\end{enumerate}
\end{proof}
 
We procede with convex separation results which are vital to the subsequent developments.

\begin{definition}
  Let 
  $C_1$ and $C_2$
  be two non-empty convex sets in $\R^n$. 
  A hyperplane $H$ is said to \textbf{separate}
  $C_1$ and $C_2$
  if $C_1$ is contained in one of the closed half-spaces associated with
  $H$ and $C_2$ lies in the opposite closed half-space. It is said to separate 
  $C_1$ and $C_2$
  \textbf{properly} if 
  $C_1$ and $C_2$
  are not \textit{both} actually contained in $H$ itselef.
\end{definition}
\begin{theorem}
  Let $C_1$ and $C_2$ be two non-empty convex sets in $\R^n$. 
  There exists a hyperplane separating
  $C_1$ and $C_2$
  properly 
  if and only if
  there exists a vector $b\in \R^n$ such that
  \begin{gather}
    \sup_{x\in C_2} \inner{x}{b}
    \le
    \inf_{x\in C_1} \inner{x}{b}
    \quad 
    \text{and}
    \quad 
    \inf_{x\in C_2} \inner{x}{b}
    <
    \sup_{x\in C_1} \inner{x}{b}
    .
  \end{gather}
\end{theorem}
\begin{proof}
  \cite[Theorem~11.1]{Rockafellar1970}
\end{proof}
\begin{ftheorem}
  \emph{(Convex separation in finite dimension)}
  Let $C_1$ and $C_2$ be two non-empty convex sets in $\R^n$. 
  Then $C_1$ and $C_2$ can be properly separated if and only if 
  $\mathrm{ri}(C_1)\cap\mathrm{ri}(C_2)=\emptyset.$
\end{ftheorem}
\begin{proof}
  \cite[Theorem~11.3]{Rockafellar1970}
\end{proof}




\begin{definition}
  Given a nonempty subset 
  $\Omega \subseteq \R^n$
  the \textbf{support function} 
  $
  \sigma_\Omega : \R^n \to \overline{\R}
  $
  of $\Omega$
  is defined by
  \begin{gather}
    \sigma_\Omega
    (x^*)
    :=
    \sup_{x \in \Omega}
    \ 
    \inner{x^*}{x}
    \qquad
    \text{for}\ 
    x^* \in \R^n
    .
  \end{gather}
\end{definition}


\begin{definition}
  Given functions
  $
    f_i:
    \R^n \to (-\infty, \infty]
  $
  for $ i = 1, \ldots, n $
  the \textbf{infimal convolution} of these functions is defined as
  \begin{gather}
    (f_1 \square \ldots \square f_m)(x)
    :=
    \inf_{
    \begin{smallmatrix}
      x_i \in \R^n \\
      \sum_{i = 1}^{m} 
        x_i
      =
      x
    \end{smallmatrix}
    }
    \sum_{i = 1}^{m}
      f_i(x_i)
  \end{gather}
\end{definition}
 
The next result establishes a connection between the support function of the intersection of two convex sets and the infimal convolution of the support functions of the sets taken by themselfes.
The proof translates the geometric concept of convex separation to the world of convex functions.

\begin{ftheorem}
  Let $C_1$ and $C_2$ be two non-empty convex sets in $\R^n$ with
  $\mathrm{ri}(C_1)\cap\mathrm{ri}(C_2)\neq\emptyset.$
  Then the support function of the intersection 
  $
    C_1 \cap C_2
  $
  is represented as
  \begin{gather}
    (\sigma_{
    C_1 \cap C_2
    })
    (x^*)
    =
    (\sigma_{C_1}\square \sigma_{C_2})
    (x^*)
    \qquad
    \text{for all}\ 
    x^* \in \R^n.
  \end{gather}
  Furthermore, for any
  $
  x^*\in \mathrm{dom}
    (\sigma_{
    C_1 \cap C_2
    })
  $
  there exist dual elements 
  $
    x_1^*
    ,
    x_2^*
    \in \R^n
  $ 
  such that 
  $
    x^*
    =
    x_1^*
    +
    x_2^*.
  $
  and
  \begin{gather}
    (\sigma_{
    C_1 \cap C_2
    })
    (x^*)
    =
    \sigma_{C_1}(x_1^*)
    +
    \sigma_{C_2}(x_2^*).
  \end{gather}
\end{ftheorem}
\begin{proof}
  \emph{\cite[Theorem~4.23]{Mordukhovich2022}}
  We want to use results on convex separation. To make the geometric property of convex separation fruitful
  to our purpose we consider two special sets. We will verify that these sets meet the requirements for convex separation, i.e., that they are convex and the intersection of their relative interiors is empty. The rest of the proof is as in the cited reference at the beginning of the proof.
  To this end, consider 
  the sets
  \begin{gather}
    \Theta_1
    \ :=\ 
    C_1 \times [\,0,\infty)
    \quad
    \text{and}
    \quad
    \Theta_2
    \ :=\ 
    \left\{ 
      (x,\lambda)\in \R^n
      \ 
      \colon
      \ 
      x \in C_2
      \ 
      \text{and}
      \ 
      \lambda
      \,
      \le
      \,
      \inner{x^*\!}{x} - \alpha
    \right\}\ .
  \end{gather}
  \todo[color=green!40,inline]{Simplify proof with properties of relative interiors.}
  Clearly, $\Theta_1$ is convex by the convexity of $C_1$. To see that $\Theta_2$ is convex consider the affine function
  $
    \varphi:
    \R^{n}\!\times \R \to \R
    ,
    \ 
    (x,\lambda)
    \mapsto
    \alpha - \inner{x^*\!}{x} - \lambda
  $ .
  From the definitions of $\varphi$ and $\Theta_2$ we get the identity
  \begin{gather*}
 \Theta_2
    \ 
    =
    \ 
    (
      C_2\!\times\R
    )
    \ 
    \cap
    \ 
    \varphi^{-1}
    (-\infty,0\,]
    \,
    .
  \end{gather*}
  Thus, by the convexity of the sets
  $C_2$
  and
  $
    \varphi^{-1}(-\infty,0]
  $
  it follows the convexity of $\Theta_2$.
  Next we show that the relative interiors of 
  $\Theta_1$ and $\Theta_2$
  do not intersect, i.e.,
  ~$
  \mathrm{ri}\,\Theta_1\cap\mathrm{ri}\,\Theta_2=\emptyset.
  $
  First note that
  \begin{gather}
    \mathrm{ri}(\Theta_1)
    =
    \mathrm{ri}(C_1)
    \times
    \mathrm{ri}([0,\infty))
    \subseteq
    \mathrm{ri}(C_1)
    \times
    (0,\infty).
  \end{gather}
  Indeed, if 
  $
    0\in
    \mathrm{ri}([0,\infty))
  $
  then there exists $t>0$ such that $-tx\ge 0$ for some $x>0.$ A contradiction.
  Furthermore
\begin{gather}
  \mathrm{ri}(\Theta_2)
  \subseteq
    \left\{ 
      (x,\lambda)\in \R^n
      \colon
      x \in \mathrm{ri}(C_2) 
      \ 
      \text{and}
      \ 
      \lambda
      <
      \inner{x^*}{x} - \alpha
    \right\}.
\end{gather}
To see this, assume there is 
$
(x,\lambda)
\in \mathrm{ri}(\Theta_2)
$
with
$
      \lambda
      =
      \inner{x^*}{x} - \alpha
      .
$
Then for some 
$
  (y,\mu)
  \in \Theta_2
$
with
$
      \mu    
      <
      \inner{x^*}{y} - \alpha
$
there exists $t>0$ such that 
$
  (x,\lambda)
  +
  t
  (
  (x,\lambda)
  -
  (y,\mu)
  )
  \in \Theta_2.
$
It follows
\begin{align}
  0
  \le
  (1+t)(\inner{x^*}{x}-\alpha -\lambda)
  +
  t
  (
    \mu 
    -\inner{x^*}{y}
    +
    \alpha
  )
  <0
,
\end{align}
a contradiction.
The first inequality is due to 
$
  (x,\lambda)
  +
  t
  (
  (x,\lambda)
  -
  (y,\mu)
  )
  \in \Theta_2
$
and the second inequality due to
$
      \mu    
      <
      \inner{x^*}{y} - \alpha
$
and
$
      \lambda
      =
      \inner{x^*}{x} - \alpha
      .
$
But then 
$
\mathrm{ri}(\Theta_1)\cap\mathrm{ri}(\Theta_2)=\emptyset
.
$
Indeed, suppose that there exists 
$
  (x,\lambda)
  \in
  \mathrm{ri}(\Theta_1)\cap\mathrm{ri}(\Theta_2)
  .
$
Then it holds
$
  \inner{x^*}{x}
  -
  \alpha
  \le
  0
$
and $\lambda>0$
since 
$
 x
 \in
  \mathrm{ri}(C_1)\cap\mathrm{ri}(C_2)
  \subseteq
  C_1\cap C_2.
$
On the other hand
\begin{gather}
  0
  <
  \lambda
  <
  \inner{x^*}{x}
  -
  \alpha
  \le
  0
  ,
\end{gather}
a contradiction.

Thus, convex separation is applicable. Confer \cite[Theorem 4.23]{Mordukhovich2022}
about the rest of the proof.
\end{proof}

\begin{takeaways}
  The support function intersection rule connects the geometric property of convex separation to an identity in terms of support function. This result is central to the analysis of convex conjugates.
\end{takeaways}

  \section{Conjugate Calculus and Fenchel-Rockafellar Theorem}
  The goal of this section is to establish the tools to calculate convex conjugates. 
We prove the conjugate sum and chain rule.
After some examples, we will derive the Fenchel-Rockafellar Theorem.
\begin{definition}
  \label{ def_convex_conjugate }
  \emph{(Convex conjugate)}
  Given a function
  $
    f:
    \R^n \to \overline{\R}
  $
  ,
  the 
  \textbf{convex conjugate}
  $
    f^*:
    \R^n \to \overline{\R}
  $
  of $f$ is defined as
  \begin{gather}
    f^*(x^*)
    :=
    \sup_{ x \in \R^n }
    (x^*)^T x - f(x)
  \end{gather}
\end{definition}

Note that $f$ in Definition~\ref{ def_convex_conjugate }
does not have to be convex. On the other hand, the convex conjugate is always convex:

\begin{proposition}
  Let  
  $
    f:
    \R^n \to ( - \infty, \infty ]
  $
  be a proper function. 
  Then its convex conjugate
  $
    f^*:
    \R^n \to ( - \infty, \infty ]
  $
  is convex.
\end{proposition}


\begin{lemma}
  For any proper function
  $
    f:\R^n\to\overline{\R}
  $
  we have
  \begin{gather}
    f^*(x^*) 
    =
    \sigma_{\mathrm{epi}(f)}
    (x^*,-1)
    \qquad
    \text{for}
    \ 
    x^* \in \R^n.
  \end{gather}
\end{lemma}
\begin{proof}
  Let $x^*\in\R^n$
  and
  $
    (x,\lambda)\in \mathrm{epi}(f).
  $
  Then
  $
    x \in \mathrm{dom}(f)
  $
  and
  $
    f(x)\le \lambda.
  $
  Thus
  \begin{gather}
    \inner{x^*}{x} - f(x)
    \ge
    \inner{x^*}{x} - \lambda
    \qquad
    \text{for all}\ 
    (x,\lambda)\in \mathrm{epi}(f).
  \end{gather}
  On the other hand 
  $
    (x,f(x))\in \mathrm{epi}(f)
  $
  for all
  $
    x \in \mathrm{dom}(f).
  $
  It follows
  \begin{gather}
    \inner{x^*}{x} - f(x)
    \le
    \sup_{(x,\lambda)\in\mathrm{epi}(f)}
    \inner{x^*}{x} - \lambda
    \qquad
    \text{for all}\ 
    x \in \mathrm{dom}(f).
  \end{gather}
  Taking the supremum in the last two displays yields
  \begin{align}
    f^*(x^*)
    =
    \sup_{x\in\mathrm{dom}(f)}
    \inner{x^*}{x} - f(x)
    &=
    \sup_{(x,\lambda)\in\mathrm{epi}(f)}
    \inner{x^*}{x} - \lambda
    \\
    &=
    \sup_{(x,\lambda)\in\mathrm{epi}(f)}
    \inner{(x^*,-1)}{(x,\lambda)} 
    =
    \sigma_{\mathrm{epi}(f)}
    (x^*,-1).
  \end{align}
\end{proof}
% conjugate chain rule %
 %%%%%%%%%%%%%%%%%%%%%%
\begin{theorem}
  \emph{(Conjugate Chain Rule)}
  \label{cvxa_conjugate_chain_rule}
  Let 
  $
    A:
      \R^m \to \R^n
  $
  be a linear map (matrix)
  and
  $
    g:
      \R^n \to (-\infty, \infty]
  $
  a proper convex function. If
  $
    \text{Im}(A) \cap \text{ri}(\text{dom}(g))
    \neq
    \emptyset
  $
  it follows
  \begin{gather}
    ( g \circ A )^* ( x^* )
    =
    \inf_
          { y^* \in ( A^* )^{ -1 } ( x^* )}
                                          g^*( y^* )
                                          .
  \end{gather}
  Furthermore, 
    for any 
      $
        x^* \in \text{dom}( g \circ A)^*
      $
        there exists
          $
            y^* \in ( A^* )^{ -1 } ( x^* )
          $
            such that
              $
                ( g \circ A)^* ( x^* )
                =
                g^*( y^* )
              $.
\end{theorem}

% conjugate sum rule %
 %%%%%%%%%%%%%%%%%%%%



\begin{theorem}
  Let
  $
    f,g:
    \R^n \to (-\infty, \infty]
  $
  be proper convex functions 
  and
  $
  \text{ri}\left( \text{dom}(f) \right)
  \cap
  \text{ri}\left( \text{dom}(g) \right)
  \neq 
  \emptyset
  .
  $
  Then we have the conjugate sum rule
  \begin{gather}
    ( f + g )^*(x^*)
    =
    ( f^* \square g^*)(x^*)
  \end{gather}
  for all $x^* \in \R^n$.
  Moreover, the infimum in 
  $
    ( f^* \square g^*)(x^*)
  $
  is attained, i.e., for any
  $
    x^* \in \text{dom}(f+g)^*
  $
  there exists vectors $x_1^*, x_2^*$
  for which
  \begin{gather}
    (f+g)^*(x^*)
    =
    f^*(x_1^*)
    +
    g^*(x_2^*),
    \quad
    x^* = x_1^* + x_2^*.
  \end{gather}
\end{theorem}
\begin{proof}
  Let $x^*\in\R^n$ and fix $x_1^*,x_2^*\in\R^n$ such that
  $x^*=x^*_1+x^*_2$.
  We get
  \begin{align*}
    f^*(x^*_1)+g^*(x^*_2)
    &=
    \sup_{x\in\R^n}
    \inner{x^*_1}{x}-f(x)
    +
    \sup_{x\in\R^n}
    \inner{x^*_2}{x}-g(x)
    \\
    &\ge
    \sup_{x\in\R^n}
    \inner{x^*_1}{x}-f(x)
    +
    \inner{x^*_2}{x}-g(x)
    =
    \sup_{x\in\R^n}
    \inner{x^*_1+x^*_2}{x}-(f(x)+g(x))
    \\
    &=
    \sup_{x\in\R^n}
    \inner{x^*}{x}-(f+g)(x)
    =(f+g)^*(x^*)
  \end{align*}
  Taking the infimum over $x_1^*,x_2^*\in\R^n$ in the above display gives 
  $
  (f^*\square g^*)(x^*)
  \ge
  (f+g)^*(x^*).
  $
  Let us prove now $\le$ under the condition
  $
  \text{ri}\left( \text{dom}(f) \right)
  \cap
  \text{ri}\left( \text{dom}(g) \right)
  \neq 
  \emptyset
  .
  $
  The only case we need to consider is
  $
    (f+g)^*(x^*)<\infty.
  $
  Define two convex sets by
  \begin{align}
    \Omega_1
    &:=
    \left\{ 
      (x,\alpha,\beta)\in\R^{n+2}
      \colon
      \alpha\ge f(x)
    \right\}
    =
    \mathrm{epi}(f)\times \R,
    \\
    \Omega_2
    &:=
    \left\{ 
      (x,\alpha,\beta)\in\R^{n+2}
      \colon
      \beta\ge g(x)
    \right\}.
  \end{align}
  Similar to Lemma we get the representation
  \begin{gather}
    (f+g)^*(x^*)
    =
    \sigma_{\Omega_1\cap\Omega_2}
    (x^*,-1,-1).
  \end{gather}
  Indeed, the only thing we need to verify is
  $
    \mathrm{dom}(f)\cap\mathrm{dom}(g)
    =
    \mathrm{dom}(f+g).
  $
  The inclusion $\subseteq$ is clear.
  Assume towards a contradiction that
  $
    (f+g)(x)<\infty
  $
  and
  $
    f(x)=\infty.
  $
  Since $g(x)>-\infty$ it holds
  \begin{gather}
    \infty
    =
    \infty+g(x)
    =f(x)+g(x)
    =(f+g)(x)
    <
    \infty.
  \end{gather}
  This is a contradiction. The same holds for $f$ and $g$ reversed. It follows the inclusion $\supseteq$ and equality.
  By the support function intersection rule there exist triples
  \begin{gather}
    (x^*_1,-\alpha_1,-\beta_1),
    (x^*_2,-\alpha_2,-\beta_2)
    \in \R^{n+2}
    \quad
    \text{such that}
    \quad
    (x^*,-1,-1)
    =
    (x^*_1+x^*_2,-(\alpha_1+\alpha_2),-(\beta_1+\beta_2))
  \end{gather}
  and
  \begin{gather}
    (f+g)^*(x^*)
    =
    \sigma_{\Omega_1\cap\Omega_2}
    (x^*,-1,-1)
    =
    \sigma_{\Omega_1}
    (x^*_1,-\alpha_1,-\beta_1)
    +
    \sigma_{\Omega_2}
    (x^*_2,-\alpha_2,-\beta_2).
  \end{gather}
  Next we show
  $\beta_1=\alpha_2=0.$
  Suppose towards a contradiction that 
  $\beta_1\neq 0.$ 
  We fix 
  $(\overline{x},\overline{\alpha})\in\mathrm{epi}(f).$
  Then
  \begin{gather}
    \sigma_{\Omega_1}
    (x^*_1,-\alpha_1,-\beta_1)
    =
    \sup_{(x,\alpha,\beta)\in \mathrm{epi}(f)\times \R}
    \inner{x^*}{x}-\alpha \alpha_1 -\beta \beta_1
    \ge
    \sup_{\beta\in \R}
    \inner{x^*}{\overline{x}}-\overline{\alpha} \alpha_1 -\beta \beta_1
    =\infty.
  \end{gather}
  This contradicts
  $
    (f+g)^*(x^*)<\infty.
  $
  In a similar fashion we can derive a contradiction for $\alpha_2\neq0.$
  Employing Lemma and taking into account the structures of the sets 
  $\Omega_1$ and $\Omega_2$ this implies
  \begin{align}
    (f+g)^*(x^*)
    &=
    \sigma_{\Omega_1\cap\Omega_2}
    (x^*,-1,-1)
    =
    \sigma_{\Omega_1}
    (x^*_1,-1,0)
    +
    \sigma_{\Omega_2}
    (x^*_2,0,-1)
    \\
    &=
    \sigma_{\mathrm{epi}(f)}(x^*_1,-1)
    +
    \sigma_{\mathrm{epi}(g)}(x^*_2,-1)
    =
    f^*(x^*_1)
    +
    g^*(x^*_2)
    \ge
    (f^*\square g^*)(x^*).
  \end{align}
  This finishes the proof.
\end{proof}




Given 
proper convex functions $f, g : \R^n \to \overline{\R}$ 
and
a matrix $A \in \R^{n \times n}$,
we define 
the primal minimization problem as follows:
\begin{problem}
  \emph{(Primal)}
  Given proper convex functions 
  $
  f:\R^n \to \overline{\R} 
  $,
  $
  g:\R^m \to \overline{\R} 
  $
  and a matrix 
  $
    A\in \R^{m \times n}
  $
  we define the \textbf{primal optimization problem} to be
  \label{cvxa_primal_problem}
  \begin{gather*}
    \underset{x \in \R^n}{\mathrm{minimize}}
    \qquad
    f(x) + g(Ax)
  \end{gather*}
\end{problem}
\begin{remark}
  Problem \autoref{cvxa_primal_problem}
  appears in the unconstrained form. We can impose constraints by controling for the domains of $f$ and $g$.
  To incorporate linear constraints $Ax \le 0$
  or more general constraints $x\in\Omega$, where $\Omega$ is a convex set,
  we can choose
  \begin{gather}
    g(x)=\delta_\Omega(x):=
  \end{gather}
  where $x\notin\Omega$
  leads to 
  $f(x) + g(x)=\infty$
  and the optimization problem (if feasible) will exclude $x$ from the solutions.
\end{remark}
\begin{problem}
  \emph{(Dual)}
  \label{cvxa_dual_problem}
  Consider the same setting as in Problem~\autoref{cvxa_primal_problem}.
  Using the convex conjugates of $f,g$ and the transpose of $A$
  we define the \textbf{dual problem} of Problem~\autoref{cvxa_primal_problem} to be
  \begin{gather*}
    \underset{y^* \in \R^m}{\mathrm{maximize}}
    \qquad
   - f^*(A^\top y^*) - g^*(y^*).
  \end{gather*}
\end{problem}
\begin{proposition}
  Consider the optimization problem~\autoref{cvxa_primal_problem} and its dual~\autoref{cvxa_dual_problem}, where the functions $f$ and $g$ are not assumed to be convex. Define the \textbf{optimal values} of these problems by
  \begin{gather*}
    \widehat{p}:= \inf_{x\in\R^n}f(x)+g(Ax)
    \quad
    \text{and}
    \quad
    \widehat{d}:= \sup_{y\in\R^m} - f^*(A^\top y) - g^*(y).
  \end{gather*}
  Then we have the relationship
  $\widehat{d}\le \widehat{p}$.
\end{proposition}
\begin{proof}
  It holds
  \begin{align*}
    - f^*(A^\top y^*) - g^*(y^*) 
    &=
      -\sup_{x\in \R^n}\inner{A^\top y^*}{x}-f(x)
      -\sup_{y\in \R^m}\inner{-y^*}{y}-g(y)
      \\
    &=
    \inf_{x\in \R^n}f(x)-\inner{y^*}{Ax}
      +\inf_{y\in \R^m}g(y)+\inner{y^*}{y}
      \\
    &\le
      \inf_{x\in \R^n}f(x)-\inner{y^*}{Ax}
      +\inf_{x\in \R^n}g(Ax)+\inner{y^*}{Ax}
      \\
    &\le
      \inf_{x\in \R^n}f(x)-\inner{y^*}{Ax} + g(Ax) + \inner{y^*}{Ax}
      \\
    &=  
      \inf_{x\in \R^n}f(x)+g(Ax)
    =
      \widehat{p}
  \end{align*}
  The first equality is due to the definition of convex conjugates, the second equality due to $\inner{A^\top y}{x}=\inner{y}{Ax}$ and $\inf \left\{ -B \right\}=-\sup \left\{ B \right\}$ for all $B \subseteq \overline{\R}$ and the first inequality due to $\mathrm{Im}(A)\subseteq \R^m$.
  Taking the supremum with respect to all 
  $y^*\in\R^m$
  yields the result.
\end{proof}
\begin{theorem}
  \label{cvxa_fenchel_theorem}
  Let 
  $f, g : \R^n \to \overline{\R}$ 
  be proper convex functions
  and
  $0 \in \text{ri}(\text{dom}(g) - A (\text{dom}(f)) )$
  .
  Then the optimal values of \eqref{cvxa_primal_problem} and \eqref{cvxa_dual_problem} are equal, 
  i.e.
  \begin{gather}
    \inf_{x \in \R^n} 
    \left\{ f(x) + g(Ax) \right\}
    =
    \sup_{y \in \R^n} \left\{   -f^* \left( A^T y \right) - g^*(-y) \right\}
    .
  \end{gather}
\end{theorem}
%%%%%%%%%%%%%%%%%%%%%%%%%%%%%%%%%%%%%%%%%%%%%%%
\begin{lemma}
  \label{syu_1_result}
  Let 
  $f : \R^n \to (-\infty, \infty]$ 
  be convex.
  Then 
  for all $y \in \R^n$ and $C>0$ 
    \begin{gather}
      \label{7060_0}
      \inf_{\norm{\Delta}=C} f(y+\Delta) - f(y) \ge 0 \quad
      \Longrightarrow
      \quad
    \exists y^* \in \R^n
    \colon
    y^* \,\text{is global minimum of $f$ and}\,
      \norm{y^* - y} \le C.
    \end{gather}
\end{lemma}
\begin{proof}
  Since 
  $\mathcal{C}:=\left\{ \norm{\Delta}\le C \right\}$
  is convex
  $f$ has a local minimum in 
  $
    y + \mathcal{C}
    :=
    \left\{ 
      y + \Delta \,
      \mid \,
      \norm{\Delta}\le C
    \right\}
    .
  $
  Suppose towards a contradiction that
  $
    y^* \in 
            y + \mathcal{C}
  $
  is a local minimum, but not a global minimum 
  and
  the left-hand side of 
  \eqref{7060_0} is true.
  Then it holds
  \begin{gather}
    \label{7060_3}
    f(x) < f(y^*)
    \quad
    \text{for some}\ 
    x 
    \in 
    \R^n 
    \setminus 
      y + \mathcal{C}
    .
  \end{gather}
  Furthermore since $y + \mathcal{C}$ is compact and contains $y^*$,
  the line segment $\mathcal{L}[y^*,x]$ contains a point on the boundary of 
  $y + \mathcal{C}$, i.e.
  \begin{gather}
    \label{7060_4}
    \theta x + (1 - \theta) y^* = y + \Delta_x
    \quad
    \text{for some}\ 
    \theta \in (0,1)\ 
    \text{and}\ 
    \Delta_x \ 
    \text{with}\ 
    \norm{\Delta_x}=C
    .
  \end{gather}
    It follows
    \begin{align}
      \label{7060_5}
      \begin{split}
      f(y^*)
      \le
      f(y)
      \le
      f(y + \Delta_x)
      &=
      f(
        \theta x + (1 - \theta) y^*
      )
      \\
      &\le
      \theta f(x)
      + 
      (1 - \theta)
      f(y^*)
      <
      f(y^*)
      ,
      \end{split}
    \end{align}
    which is a contradiction.
    Thus every local minimum of $f$ in $y + \mathcal{C}$ is also a global minimum.
    The first inequality is due to
    $y^*$ being a local minimum of $f$ in
    $
      y + \mathcal{C},
    $
    the second inequality is due to the left-hand side of 
    \eqref{7060_0} being true,
    the equality is due to \eqref{7060_4},
    the third inequality is due to the convexity of $f$
    and the strict inequality is due to \eqref{7060_3}.
    %It follows the right-hand side of \eqref{7060_0}.
\end{proof}


\begin{takeaways}
  Almost there 
  \lipsum[1]
\end{takeaways}

%%%%%%%%%%%%%%%%%
% INTRODUCTION %
%%%%%%%%%%%%%%%%
In the following we do not expect the reader to be familiar with convex analysis. However, some very well known results will be stated without proof. The interested reader can study \cite{Mordukhovich2022} for the bedrock analysis.


We begin by defining convex sets
%

\begin{definition}
  A subset $\Omega\subseteq \R^n$ is called CONVEX if we have $\lambda x+(1-\lambda)y\in \Omega$ for all $x,y\in \Omega$ and $\lambda\in (0,1)$. 
\end{definition}

Clearly, the line segment 
$[a,b]:=\left\{ \lambda a+(1-\lambda)b\,\mid \, \lambda\in [0,1] \right\}$ is contained in $\Omega$ for all $a,b\in \Omega$ if and only if $\Omega$ is a convex set.
%

Next we define convex functions. 
%

The concept of convex functions is closely related to convex sets.
%  
 
The line segment between two points on the graph of a convex function lies on or above and does not intersect the graph.
%

In other words: The area above the graph of a convex function $f$ is a convex set, i.e. the \textit{epigraph}
$\text{epi}(f):=\left\{ (x,\alpha)\in \R^n\times\R\,\mid\, f(x)\le \alpha\right\}$ is a convex set in $\R^{n+1}$.
%

Often an equivalent characterisation of convex functions is more useful.
%

\begin{theorem}
  The convexity of a function $f:\R^n\to \overline{\R}$ on $\R^n$ is equivalent to the following statement:

  For all $x,y\in \R^n$ and $\lambda\in(0,1)$ we have 
    \begin{align}
      f(\lambda x + (1-\lambda)y)\le \lambda f(x)+(1-\lambda)f(y).
    \end{align}
\end{theorem}


\begin{definition}
  \emph{proper convex function}
\end{definition}

\begin{definition}
  \label{ def_convex_conjugate }
  \emph{(Convex Conjugate)}
  Given a function
  $
    f:
    \R^n \to ( - \infty, \infty ]
  $
  ,
  the convex conjugate
  $
    f^*:
    \R^n \to ( - \infty, \infty ]
  $
  of $f$ is defined as
  \begin{gather}
    f^*(x^*)
    :=
    \sup_{ x \in \R^n }
    (x^*)^T x - f(x)
  \end{gather}
\end{definition}

Note that $f$ in Definition~\ref{ def_convex_conjugate }
does not have to be convex. On the other hand, the convex conjugate is always convex:

\begin{proposition}
  Let  
  $
    f:
    \R^n \to ( - \infty, \infty ]
  $
  be a proper function. 
  Then its convex conjugate
  $
    f^*:
    \R^n \to ( - \infty, \infty ]
  $
  is convex.
\end{proposition}

%%%%%%%%%%%%%%%%%%%%%%
% CONJUGATE CALCULUS %
%%%%%%%%%%%%%%%%%%%%%%
\section{Conjugate Calculus}
When studying different primal problems such as \eqref{primal_weighting_binary} we often turn to the dual instead.
Therefore we need some reliable tools.
Beign able to compute specific convex conjugates is one tool required.



%%%%%%%%%%%%%%%%%%%
% FENCHEL DUALITY %
%%%%%%%%%%%%%%%%%%%
\section{Fenchel Duality}


Given 
proper convex functions $f, g : \R^n \to \overline{\R}$ 
and
a matrix $A \in \R^{n \times n}$,
we define 
the primal minimization problem as follows:

\begin{gather}
  \label{cvxa_primal_problem}
  \text{minimize}
  \quad
  f(x) + g(Ax) 
  \quad
  \text{subject to}
  \quad
  x \in \R^n.
\end{gather}

The Fenchel dual problem is then

\begin{gather}
  \label{cvxa_dual_problem}
  \text{maximize}
  \quad
  -f^* \left( A^T y \right) - g^*(-y) 
  \quad
  \text{subject to}
  \quad
  y \in \R^n.
\end{gather}

\begin{theorem}
  \label{cvxa_fenchel_theorem}
  Let 
  $f, g : \R^n \to \overline{\R}$ 
  be proper convex functions
  and
  $0 \in \text{ri}(\text{dom}(g) - A (\text{dom}(f)) )$
  .
  Then the optimal values of \eqref{cvxa_primal_problem} and \eqref{cvxa_dual_problem} are equal, 
  i.e.
  \begin{gather}
    \inf_{x \in \R^n} 
    \left\{ f(x) + g(Ax) \right\}
    =
    \sup_{y \in \R^n} \left\{   -f^* \left( A^T y \right) - g^*(-y) \right\}
    .
  \end{gather}
\end{theorem}

  \section{Tseng Bertsekas}
  We present the relevant parts of the paper \cite{Bertsekas2003}.

Consider the following optimization problem

  \begin{gather*}
    \underset{x\in \R^m}{\mathrm{minimize}}
    \qquad
    f(x)
  \end{gather*}
subject to the constraints
\begin{gather}
  \mathbf{A}x \ge b
  ,
\end{gather}
Where 
$
  f:
  \R^m
  \to 
  \overline{\R}
  ,
  \ 
  \mathbf{A} \ 
  \text{is a given}\ 
  n\times m
  \ 
  \text{matrix,}
$
 and $b$ is a vector in $\R^n.$
\todo[color=green!40, inline]{
Generalize also to take equality constraints.
Write in unconstrained form to derive dual.
}
\begin{assumption}
  \begin{enumerate}[label={(\roman*)}]
    Assume that the map 
    $
      f: \R^m \to \overline{\R}
    $
    has the following properties.
    \item
      $
        f 
        \ 
        \text{is strictly convex.}
      $
    \item
      $
        f
        \ 
        \text{is lower-semicontinuous and continuous}
        \ 
        \mathrm{dom}(f)
        .
      $
    \item
      $
        \text{The convex conjugate}
        \ 
        f^*
        \ 
        \text{of}\ 
        f
        \ 
        \text{is finite}
        .
      $
  \end{enumerate}
\end{assumption}

The dual optimization problem associated with $(P)$
is 


  \begin{gather*}
    \underset{p\in \R^n}{\mathrm{maximize}}
    \qquad
    q(p)
  \end{gather*}
subject to the constraints
\begin{gather}
  p \ge 0
  ,
\end{gather}
where 
$
 q: \R^n \to \overline{\R}
$
is the concave function given by
\begin{gather}
  q(p)
  :=
  \min_{x \in \R^m}
  f(x)
  +
  \inner{p}{b - \mathbf{A}x}
  =
  \inner{p}{b}
  -
  f^*(\mathbf{A}^\top p)
  .
\end{gather}
The dual problem $(D)$ is a concave program with simple nonnegativity constraints. Furthermore, strong duality holds for $(P)$ and $(D)$
, i.e., the optimal value of $(P)$
equals the optimal value of $(D).$

Since $f^*$ is real-valued and $f$ is strictly convex, 
$f^*$ and $q$ are continuously differentiable.

\begin{theorem}
  \emph{\cite[Theorem~26.3]{Rockafellar1970}}
  A closed proper convex function is (essentially) strictly convex
  if and only if its conjugate is essentially smooth.
\end{theorem}
\todo[color=red!40,inline]{Read and understand proof (p.270) }

We will denote the gradient of $q$ at $p$ by $d(p)$
and its $i$th coordinate by $d_i(p).$
Since $q$ is continuously differentiable, $d_i(p)$ is continuous, 
and since $q$ is concave, 
$d_i(p)$ as nonincreasing in $p_i.$

By differentiating and by using the chain rule, 
we obtain the dual cost gradient
\begin{gather}
  d(p)=b-\mathbf{A}x,
  \quad
  \text{where}
  \quad
  x:=
  \nabla
  f^*
  (
    \mathbf{A}^\top
    p
  )
  =
  \mathrm{argsup}
  _{\xi \in \R^m}
  \inner{p}{\mathbf{A}\xi}
  - f(\xi)
  .
\end{gather}
The last equality follows from Danskin's Theorem and \cite[Theorem~23.5]{Rockafellar1970}


\todo[color=red!40,inline]{Read and understand proof (p.80) }
\begin{proposition}
  \emph{(Danskin's Theorem \cite[page 649]{Bertsekas2003})}
  Let 
  $
    Z \subseteq \R^m
  $
  be a non-empty set, 
  and let 
  $
    \phi :
    \R^n \times Z \to \R
  $
  be a continuous function such that
  $
    \phi(\cdot,z)
    :
    \R^n \to \R,
  $
  viewed as a function of its first argument, 
  is convex for each 
  $
    z \in Z.
  $
  Then the function 
  \begin{gather}
    f:\R^n \to \R
    ,
    \qquad
    x\mapsto
    \sup_{z\in Z}
    \phi(x,z)
  \end{gather}
  is convex and has directional derivative given by
  \begin{gather}
    f^{'}(x;y)
    =
    \sup_{z \in Z(x)}
    \phi^{'}(x,z; y)
    ,
  \end{gather}
  where 
  $
    \phi^{'}(x,z; y)
  $
  is the directional derivative of the function 
  $
    \phi(\cdot,z)
  $
  at
  $x$ in the direction $y,$
  and
  \begin{gather}
    Z(x)
    :=
    \left\{ 
      \overline{z} \in \R^m
      \colon
      \phi(x, \overline{z})
      =
    \sup_{z\in Z}
    \phi(x,z)
  \right\}
  .
  \end{gather}
  In particular, if $Z(x)$ consists of a unique point 
  $\overline{z}$
  and
  $
    \phi(\cdot, \overline{z})
  $
  is differentiable at $x,$
  and 
  $
    \nabla f(x)
    =
    \nabla_x
    \phi(x,\overline{z})
    ,
  $
  where 
  $
    \nabla_x
    \phi(x,\overline{z})
  $
  is the vector with coordinates
  $
    (\partial \phi / \partial x_i)
    (x,\overline{z})
  $
\end{proposition}

Note that $x$ is the unique vector satisfying 
\begin{gather}
  \mathbf{A}p
  \in 
  \partial
  f(x)
  .
\end{gather}
From the optimality conditions for $(D)$ it follows that 
a dual vector is an optimal solution of $(D)$
if and only if 
\begin{gather}
  p
  =
  [p + d(p)]^+
  ,
\end{gather}
where $[\ \cdot \ ]^+$ is the projection onto the positive orthant, i.e.,
$
  [y]^+
  =
  [
    0 \lor y_1,
    \ldots
    0 \lor y_n,
  ]
  ^\top
  .
$
\todo[color=orange!40,inline]{Provide details. See notes.}

Given an optimal dual solution $p$,
we may obtain an optimal primal solution from the equation 
$
  x=
  \nabla
  f^*
  (
    \mathbf{A}^\top
    p
  )
.
$
To see this, 
note that 
 \begin{gather}
   \mathbf{A}x \ge b
   \quad
   \text{and}
   \quad
   p_i 
   =
   0
   \quad 
   \text{for all}
   \ 
   i
   \ 
   \text{such that}
   \quad 
   \sum_{j=1}^{m} 
   a_{ij}
   x_j
   >
   b_i
   .
 \end{gather}

 We can show that $p$ and $x$ satisfy the KKT conditions and thus $x$ is an optimal solution to $(P)$.

n 

\begin{definition}
  \emph{\cite[§28]{Rockafellar1970}}
  By an \textbf{ordinary convex program}
  $(P)$
  we mean an optimization problem of the following form
  \begin{gather*}
    \underset{x\in C}{\mathrm{minimize}}
    \qquad
    f_0(x)
  \end{gather*}
subject to the constraints
\begin{gather}
  f_1(x)\le 0,
  \ldots,
  f_r(x)\le 0,
  \qquad
  f_{r+1}(x)= 0,
  \ldots,
  f_m(x)= 0,
\end{gather}
where $C\subseteq \R^n$
is a non-empty convex set,
$f_i$ is a finite convex function on $C$ for $i\in \left\{ 1,\ldots,r \right\}$
 and 
 $f_i$ is an affine function on $C$ for $i\in \left\{ r+1, \ldots, m \right\}.$
\end{definition}

\begin{definition}
  We define 
  $
    [\lambda_1, \ldots, \lambda_m]\in \R^m
  $
  to be a \textbf{Karush-Kuhn-Tucker (KKT) vector}
  for $(P)$, if
  \begin{enumerate}[label={(\roman*)}]
    \item
      $
        \lambda_i \ge 0
        \ 
        \text{for all}
        \ 
        i\in \left\{ 1,\ldots,r \right\}.
      $
    \item
      The infimum of the proper convex function 
      $
        f_0
        +
        \sum_{i=1}^{m}
        \lambda_1 f_i
      $
      is finite and equal to the optimal value in $(P).$
  \end{enumerate}
\end{definition}

\begin{ftheorem}
  \emph{(Karush-Kuhn-Tucker conditions)}
  Let $(P)$
  be an ordinary convex program,
  $
  \overline{\alpha}
  \in \R^m
  $,
   and 
   $
   \overline{z}
   \in \R^n.
   $
   Then 
   $
  \overline{\alpha}
   $
   is a KKT vector for $(P)$
   and 
   $
   \overline{z}
   $
   is an optimal solution to $(P)$
   if and only if 
   $
   \overline{z}
   $
   and 
   the components $\alpha_i$ of
   $
  \overline{\alpha}
   $
   satisfy 
   the following conditions.

  \begin{enumerate}[label={(\roman*)}]
    \item
      $
        \alpha_i \ge 0,
        \ 
        f_i(
   \overline{z}
        )
        \le 0,
        \ 
        \text{and}
        \ 
        \alpha_i 
        f_i(
   \overline{z}
        )
        =0
        \ 
        \text{for all}
        \ 
        i\in \left\{ 1, \ldots, r \right\}
        .
      $
      \item
        $
        f_i(
   \overline{z}
        )
        =0
        \ 
        \text{for}
        \ 
        i\in \left\{ r+1, \ldots, m \right\}
        .
        $
      \item
        $
         0
         _n
         \in 
         [
          \partial
        f_0(
   \overline{z}
        )
        +
        \sum_{\alpha_i\neq 0}
        \alpha_i 
        \partial
        f_i(
   \overline{z}
        )
         ]
.
        $
  \end{enumerate}
\end{ftheorem}
\begin{proof}
  \cite[Theorem~28.3]{Rockafellar1970}
\end{proof}

\begin{takeaways}
  The Karush-Kuhn-Tucker conditions allow us
  to derive duality in terms of the optimal solutions
  for strictly convex functions.
\end{takeaways}



\chapter{Random Matrix Inequalities}
  In our application we want to bound moments of vector-valued random variables.
  For this we choose the theory of random matrix inequalities
  which lately received a lot of attention.
  In particular an approach via the method of exchangable pairs \cite{Mackey2014}
  has been fruitful in simplifying the proofs of long standing results such as the matrix Khintchin inequality.
  We base our exposition on \cite{Mackey2014}. A lot will be exact copy of this paper, so no originality is claimed. Where it seemed fit, we conducted some calculations in more detail than presented in the paper. 

  We will first introduce the method of exchangable pairs and derive auxiliary theorems to establish the matrix Khintchin inequality.
  Then we will derive inequalities for moments of matrices, first for psd matrices and then via the Hermitian dilataition for general rectangular matrices. In a last step we will introduce the notion of intrinsic dimension to improve the bounds.

  \section{A Matrix Analysis Primer}
    The \textbf{trace} of a square matrix, denoted by $\mathrm{tr},$
  is the sum of its diagonal entries, i.e. 
  $
    \mathrm{tr}(\mathbf{B})
    =
    \sum_{j=1}^{d}b_{jj}
    \quad 
    \text{for}\ 
    \mathbf{B} \in \mathbb{M}_d.
  $
  The trace is unitarily invariant, i.e.
  $
    \mathrm{tr}(\mathbf{B})
    =
    \mathrm{tr}(\mathbf{Q}\mathbf{B}\mathbf{Q}^*)
    \quad 
    \text{for all}
    \ 
    \mathbf{B}\in \mathbb{M}_d
    \ 
    \text{for all unitary}\ 
    \mathbf{Q} \in \mathbb{M}_d.
  $
  In particular, the existence of an eigenvalue value decomposition shows 
  that the trace of a Hermitian matrix equals the sum of its  eigenvalues.
  Let
  $
  f: I\to \R
  $
  where 
  $I\subseteq\R$ 
  is an interval.
  Consider a matrix 
  $\mathbf{A}\in \mathbb{H}_d$
  whose eigenvalues are contained in $I.$
  We define the matrix 
  $
    f(\mathbf{A})\in \mathbb{H}_d
  $
  using an eigenvalue decomposition of $\mathbf{A}:$
  \begin{gather}
    f(\mathbf{A})
    =
    \mathbf{Q}
    \begin{bmatrix}
      f(\lambda_1) &&\\
                   &\ddots&\\
                   && f(\lambda_d)
    \end{bmatrix}
    \mathbf{Q}^*
    \qquad
    \text{where}
    \qquad
    \mathbf{A}
    =
    \mathbf{Q}
    \begin{bmatrix}
      \lambda_1 &&\\
                   &\ddots&\\
                   && \lambda_d
    \end{bmatrix}
    \mathbf{Q}^*
    =
    \sum_{i=1}^{d} 
    \lambda_i
    \mathbf{Q}_{\bullet i}
    \mathbf{Q}_{\bullet i}^*
    .
  \end{gather}
  The definition of $f(\mathbf{A})$ does not depend on which 
  eigenvalue decomposition we choose.
  Any matrix function that arises in this fashion is called a \textbf{standard matrix function}.

  For each $p\ge 1$
  the \textbf{Schatten $p$-norm} is defined as 
  $
    \norm{\mathbf{B}}_p
    :=
    (
    \mathrm{tr}(\left| \mathbf{B} \right|^p)
    )
    ^{1/p}
    \ 
    \text{for}
    \ 
    \mathbf{B} \in \mathbb{M}_d.
  $
  In this setting,
  $
  \left| \mathbf{B} \right|
  :=
  (
    \mathbf{B}^*
    \mathbf{B}
    )^{1/2}
    .
  $
  The \textbf{spectral norm}
  of an Hermitian matrix $\mathbf{A}$
  is defined by the relation
  $
    \norm{\mathbf{A}}
    :=
    \lambda_{\mathrm{max}}(\mathbf{A})
    \lor
    (
    -\lambda_{\mathrm{min}}(\mathbf{B})
    )
    .
  $
  For a general matrix $\mathbf{B},$
  the spectral norm is defined to be the largest singular value:
  $
    \norm{\mathbf{B}}
    :=
    \sigma_1(\mathbf{B})
    .
  $
  The Schatten $p$-norm dominates the spectral norm for all $p\ge 1.$



\begin{proposition}
  Let
  $
  f,g: I\to \R
  $
  be real-valued functions on an interval $I\subseteq\R,$ 
  and let
  $\mathbf{A}\in \mathbb{H}_d$
  be a Hermitian matrix
  whose eigenvalues are contained in $I.$

  \begin{enumerate}[label={(\roman*)}]
    \item
      If $\lambda$ is an eigenvalue of of $\mathbf{A},$
      then $f(\lambda)$ is an eigenvalue of $f(\mathbf{A}).$
    \item
      $
        f(a)
        \le
        g(a)
        \quad
        \text{for all}\ 
        a\in I
        \quad
        \text{implies}
        \quad
        f(\mathbf{A})
        \preccurlyeq
        g(\mathbf{A})
        .
      $
  \end{enumerate}
\end{proposition}


\begin{takeaways}
  This Primer is not a prim number.
  \lipsum[5]
\end{takeaways}

  \section{The Method of Exchangeable Pairs}
  \subsection{Matrix Stein pairs}
We first define an exchangable pair.
\begin{definition}
  Let 
  $Z$ and $Z^{'}$
  random variables taking values
  in a Polish space $\mathcal{Z}.$
  We say that
  $
  (
    Z
    ,
    Z^{'}
  )
  $
  is an \textbf{exchangable pair}
  if it has the same distribution as 
  $
  (
  Z^{'}
    ,
    Z
  )
  .
  $
  In particular, 
  $Z$ and $Z^{'}$
  must share the same distribution.
\end{definition}

The following approach originates in the work of Charles Stein \cite{Stein1972} 
on normal approximation for a sum of dependent random variable.
We will explain how some central ideas of this theory extends to matrices.

We can obtain a lot of information about the fluctuation of a random matrix
$\mathbf{X}$
if we can construct a good exchangable pair 
$
(
  \mathbf{X}
  ,
  \mathbf{X}^{'}
)
.
$
With this motivation in mind, let us introduce a special class of exchangable pairs.
\begin{definition}
  Let
  $
  (
    Z
    ,
    Z^{'}
  )
  $
  be an exchangable pair of random variables taking values
  in a Polish space $\mathcal{Z},$
  and let 
  $
    \mathbf{\Psi}
    : 
    \mathcal{Z}
    \to 
  \mathbb{H}_d
  $
  be a measurable function.
  Define the random Hermitian matrices
  \begin{gather}
    \mathbf{X}
    :=
    \mathbf{\Psi}
    (Z)
    \quad
    \text{and}
    \quad
    \mathbf{X}^{'}
    :=
    \mathbf{\Psi}
    (Z^{'})
    .
  \end{gather}
  We say that 
  $
  (
    \mathbf{X}
    ,
    \mathbf{X}^{'}
  )
  $
  is a \textbf{matrix Stein pair}
  if there is a constant 
  $\alpha\in (0,1]$
  for which
  \begin{gather}
    \E[
    \mathbf{X}
    -
    \mathbf{X}^{'}
    \vert
    Z
    ]
    =
    \alpha 
    \mathbf{X}
    \qquad
    \text{almost surely.}
  \end{gather}
  The constant 
  $\alpha$
  is called the \textbf{scale factor} of the pair.
  We always assume 
  $
    \E
    \left[ 
      \norm{\mathbf{X}}^2
    \right]
    <
    \infty
    .
  $
\end{definition}

A matrix Stein pair 
$
(
\mathbf{X}
,
\mathbf{X}^{'}
)
$
has several useful propreties. First,
$
(
\mathbf{X}
,
\mathbf{X}^{'}
)
$
always forms an exchangable pair. Second, it must be the case that
$\E[\mathbf{X}]=\mathbf{0}.$
Indeed,
\begin{gather*}
  \E[\mathbf{X}]
  =
  \frac{1}{\alpha}
  \E[
  \E
  [
    \mathbf{X}
    -
    \mathbf{X}^{'}
    \vert
    Z
  ]
  ]
  =
  \frac{1}{\alpha}
  \E[
    \mathbf{X}
    -
    \mathbf{X}^{'}
    \vert
  ]
  =
  \mathbf{0}
  .
\end{gather*}
\subsection{The method of exchangable pairs}
A well-chosen matrix Stein pair 
$
(
\mathbf{X}
,
\mathbf{X}^{'}
)
$
provides a surprisingly powerful tool for studying 
the random matrix $\mathbf{X}.$
The technique depends on a fundamental technical lemma.
\begin{lemma}
  Suppose that
  $
  (
    \mathbf{X}
    ,
    \mathbf{X}^{'}
  )
  $
  is a matrix Stein pair with scale factor $\alpha.$
  Let 
  $
    \mathbf{F}
    :
    \mathbb{H}_d
    \to
    \mathbb{H}_d
  $
  be a measurable function that satisfies the regularity condition
  $
  \E
  \left[
  \norm{
    (
    \mathbf{X}
    -
    \mathbf{X}^{'}
  )
  \mathbf{F}(\mathbf{X})
  }
  \right]
  <
  \infty
  .
  $
  Then
\begin{gather}
  \E
  [
    \mathbf{X}
    \cdot
    \mathbf{F}
    (\mathbf{X})
  ]
  =
  \frac{1}{2\alpha}
  \E
  [
    (
    \mathbf{X}
    -
    \mathbf{X}^{'}
    )
    (
    \mathbf{F}
    (
    \mathbf{X}
    )
    -
    \mathbf{F}
    (
    \mathbf{X}^{'}
    )
  )
  ]
  .
\end{gather}
\end{lemma}
In short, the randomness in the Stein pair furnishes an alternative expression for the expected product of $\mathbf{X}$
and a function $\mathbf{F}.$
It allows us to estimate the expectation using the smoothness properties of the function $\mathbf{F}$ and the discrepancy between $\mathbf{X}$
and $\mathbf{X}^{'}.$
\begin{proof}
  \emph{\cite[Lemma~2.4]{Mackey2014}}
  Suppose that
  $
  (
    \mathbf{X}
    ,
    \mathbf{X}^{'}
  )
  $
  constructed from an auxiliary exchangable pair 
  $
  (
  Z,
  Z^{'}
  )
  .
  $
  The defining property implies
 \begin{gather}
   \alpha \cdot 
   \E[
   \mathbf{X}
   \cdot
   \mathbf{F}
   (\mathbf{X})
   ]
   =
   \E[
    \E[
    \mathbf{X}
    -
    \mathbf{X}^{'}
    \vert
    Z
    ]
    \cdot
   \mathbf{F}
   (\mathbf{X})
   ]
   =
    \E[
    (
    \mathbf{X}
    -
    \mathbf{X}^{'}
    )
   \mathbf{F}
   (\mathbf{X})
   ]
 \end{gather} 
\end{proof}
\subsection{The conditional variance}
To each matrix Stein pair 
  $
  (
    \mathbf{X}
    ,
    \mathbf{X}^{'}
  ),
  $
  we may associate a random matrix called the conditional variance of $\mathbf{X}.$
  The purpose of this section is to argue that the spectral norm of $\mathbf{X}$
  is unlikely to be large, when the conditional variance is small.
\begin{definition}
  Suppose that
  $
  (
    \mathbf{X}
    ,
    \mathbf{X}^{'}
  ),
  $
  is a matrix Stein pair, constructed from an auxiliary exchangeable pair
  $
  (
  Z
    ,
  Z^{'}
  ).
  $
  The \textbf{conditional variance}
  is the random matrix
  \begin{gather}
    \mathbf{\Delta_X}
    :=
    \mathbf{\Delta_X}
    (Z)
    :=
    \frac{1}{2\alpha}
    \E
    [
    (
    \mathbf{X}
    -
    \mathbf{X}^{'}
    )^2
    \vert
    Z
    ]
    ,
  \end{gather}
  where $\alpha$ is the scale factor of the pair. We may take any version of the conditional expectation in this definition.
\end{definition}

The conditional variance
$
    \mathbf{\Delta_X}
$
can be regarded as a stochastic estimate for the variance 
of the random matrix $\mathbf{X}.$
To see this, assume
the second moment of $\mathbf{X}$ exists. Then it follows from Lemma 
with $\mathbf{F}(\mathbf{X})=\mathbf{X}$
\begin{gather}
  \E
  [
    \mathbf{\Delta_X}
  ]
  =
  \E
  [
  \mathbf{X}^2
  ]
  .
\end{gather}
To verify the regularity condition, note that
\begin{gather}
  \E
  [
  \lVert
    (
    \mathbf{X}
    -
    \mathbf{X}^{'}
    )
    \mathbf{X}
    \rVert
  ]
  \le
  \E[
  \norm{
    \mathbf{X}
  }
  ^2
  ]
  +
  \E
  [
  \norm{
    \mathbf{X}
  }
  \cdot
  \lVert
    \mathbf{X}^{'}
    \rVert
  ]
  \le
  2
  \E
  [
  \norm{
    \mathbf{X}
  }
  ^2
  ]
  <
  \infty
  .
\end{gather}
\begin{example}
  \emph{\cite[Example~2.4]{Mackey2014}}
\end{example}
nrt\\
is not necessary, as, for example, the normal distribution does not fulfill it. For

 \section{Matrix Khintchin Inequality and Applications}
 The goal of this section is to derive the matrix Khintchin inequality and show some important applications.
 For this we need an auxiliary theorem which is an extension of the 
 \textit{Burkholder-Davis-Gundy (BDG) inequality} from classical martingale theory~\cite{Burkholder1973}.
  We prepare for the proof of this theorem by assembling some analytic tools.
  \begin{theorem}
  \emph{(Matrix BDG inequality)}
  Let
  $
    p = 1
    \ 
    \text{or}\ 
    p \ge 3/2
    .
  $
  Suppose that 
  $
  (
    \mathbf{X}
   , 
   \mathbf{X}^{'}
  )
  $
  is a matrix Stein pair where
  $
   \E
   [ 
    \norm{\mathbf{X}}
    _{2p}^{2p}
   ]
   <
   \intfy
   .
  $
\end{theorem}




\begin{ftheorem}
  \emph{\cite[Corollary~7.3]{Mackey2014}}
  Suppose that
  $
    p = 1
    \ 
    \text{or}\ 
    p \ge 3/2
    .
  $
  Consider a finite sequence
  $
    (\mathbf{Y}_k)_{k\ge 1}
  $
  of independent, random, Hermitian matrices 
  and a deterministic sequence
  $
    (\mathbf{A}_k)_{k\ge 1}
  $
  for which
  \begin{gather}
    \E[\mathbf{Y}_k]
    =
    0
    \quad 
    \text{and}
    \quad
    \mathbf{Y}_k^2
    \preccurlyeq
    \mathbf{A}_k^2
    \qquad
    \text{almost surely for all}\ 
    k \ge 1.
  \end{gather}
  Then
  \begin{gather}
      \E
      \left[
        \norm{
          \sum_{k\ge 1}
            \mathbf{Y}_k
        }
        _{2p}
        ^{2p}
      \right]
      ^{1/(2p)}
      \le
      \sqrt{
        p - \frac{1}{2}
      }
      \,
      \norm{
        \left( 
          \sum_{k\ge 1}
          (
            \mathbf{A}_k^2
            + 
            \E[
              \mathbf{Y}_k^2
            ]
          )
        \right)
        ^{1/2}
        }
      _{2p}
      .
  \end{gather}
  In particular, when 
  $
    (\xi_k)_{k\ge 1}
  $
  is an independent sequence of Rademacher random variables,
  \begin{gather}
      \E
      \left[
        \norm{
          \sum_{k\ge 1}
            \xi_k
            \mathbf{A}_k
        }
        _{2p}
        ^{2p}
      \right]
      ^{1/(2p)}
      \le
      \sqrt{
        2p - 1
        }
      \,
      \norm{
        \left( 
          \sum_{k\ge 1}
            \mathbf{A}_k^2
        \right)
        ^{1/2}
        }
      _{2p}
      .
  \end{gather}
\end{ftheorem}

  \section{Generalzed Inequalities by Hermitian Dilataition}
  For an introduction to Hermitian Dilation see \cite[§2.1.16]{Tropp2015}
\begin{definition}
  \emph{(Hermitian Dilation)}
  The Hermitian dilation
  \begin{gather*}
    \mathfrak{H} : \C^{d_1 \times d_2} \to \mathbb{H}_{d_1 \times d_2}
  \end{gather*}
  is a map from a general matrix to an Hermitian matrix defined by
  \begin{gather}
    \label{ rmineq_hermitian_dilation } 
    \mathfrak{H}(B)
    :=
    \begin{bmatrix}
      0   & B \\
      B^* & 0 \\
    \end{bmatrix}
  \end{gather}
\end{definition}

Properties:

$
\norm{\mathbf{B}}
=
\norm{\mathfrak{H}\mathbf{B}}
$
Linearity
confer \cite[§7.7.3]{Tropp2015}
for combination with intrinsic dimension argument.

  \begin{ftheorem}
  %\label{rmineq_bernstein}
  Let $\mathbf{A}_1, \ldots, \mathbf{A}_n$ be independent, random matrices with dimension 
  $d_1 \times d_2$.
    Introduce the random matrix
      \begin{gather*}
        \mathbf{S}:=\sum_{k=1}^n \mathbf{A}_k.
      \end{gather*}
    Let $v(\mathbf{S})$ be the matrix variance statistic of the sum:
      \begin{align}
        \label{rmineq_bernstein_cond_2}
        v(\mathbf{S}):= \norm{\E[\mathbf{S}\mathbf{S}^\top ]} \lor \norm{\E[\mathbf{S}^\top \mathbf{S}]} 
             = \norm{\sum_{k=1}^n\E[\mathbf{A}_k\mathbf{A}_k^\top]} \lor \norm{\sum_{k=1}^n\E[\mathbf{A}_k^T \mathbf{A}_k]} .
      \end{align}
    Then
      \begin{align}
        \label{rmineq_bernstein_expectation_bound}
        \left(
          \E \left[ \norm{\mathbf{S}}^2 \right]
        \right)^{\frac{1}{2}}
        \le
        \sqrt{
          2ev(\mathbf{S})\log(d_1+d_2)
        } 
        + 
        4e \left( 
          \E[\max_{k \le n}\norm{\mathbf{A}_k}^2]
        \right)^\frac{1}{2}
        \log(d_1+d_2).
      \end{align}
\end{ftheorem}
  

  \begin{takeaways}
    Dilataition is so deep.
    \lipsum[3]
  \end{takeaways}
  
  \section{Intrinsic Dimension}
  \begin{definition*}
  \label{rmineq_intrinsic_bernstein}
  For a positive-semidefinite matrix $\,\mathbf{S}$,
  the \textbf{intrinic dimension} is the quantity
  \begin{gather*}
    \mathrm{intdim}
    \,
    \mathbf{A}
    \ 
    :=
    \ 
\mathrm{tr}\,\mathbf{A}
\,
/
\,
\norm{\mathbf{A}}
    \,
    .
  \end{gather*}
\end{definition*}
\begin{lemma*}
  \emph{(Intrinsic dimenision)}
  Let 
  $
  \,
    \varphi\,:\, [\,0,\infty) \,\to\, \R
    \,
  $
  be a convex function with
  $
    \varphi(0)\,=\,0
  $.
  For any positive-semidefinite matrix $\,\mathbf{S}$ it holds 
  \begin{gather*}
    \mathrm{tr}\,\varphi(\mathbf{S})
    \ 
    \le
    \ 
    \mathrm{intdim}\,\mathbf{S}
    \ 
    \cdot
    \ 
    \varphi\!\norm{\mathbf{S}}
    \,
    .
  \end{gather*}
\end{lemma*}
\begin{proof}
  \emph{\cite[Lemma~7.5.1]{Tropp2015}}
  Since $\varphi$ is convex on any interval $[\,0,L\,]$ with $L\,>\,0$, and $\varphi(0)\,=\,0$, it holds
  \begin{gather*}
    \varphi(a)
    \ 
    \le
    \ 
    \left( 
      1 - a/L
    \right)
    \cdot
    \varphi(0)
    \ 
    +
    \ 
    a/L
    \cdot
    \varphi(L)
    \ 
    =
    \ 
    a/L
    \cdot
    \varphi(L)
    \qquad
    \text{for all}\ 
    a \,\in\, [\,0,L\,]
    \,
    .
  \end{gather*}
  Since $\,\mathbf{S}\,$ is positive-semidefinite, the eigenvalues of $\,\mathbf{S}\,$ 
  fall in the interval $\,[\,0,L\,]\,$, where $\,L\,=\,\norm{\mathbf{S}}\,$.
  It follows 
  \begin{align*}
    \mathrm{tr}\,\varphi(\mathbf{S})
    &
    \ 
    =
    \ 
    \sum_{i=1}^{d}
    \varphi(\lambda_i)
    \ 
    \le
    \ 
    \sum_{i=1}^{d}
    \lambda_i
    /
    \norm{\mathbf{S}}
    \ 
    \cdot
    \ 
    \varphi(\norm{\mathbf{S}})
    \\
    &
    \ 
    =
    \ 
    \mathrm{tr}(\mathbf{S})
    /
    \norm{\mathbf{S}}
    \ 
    \cdot
    \ 
    \varphi(\norm{\mathbf{S}})
    \ 
    =
    \ 
    \mathrm{intdim}\,\mathbf{S}
    \ 
    \cdot
    \ 
    \varphi\!\norm{\mathbf{S}}
    \,
    .
  \end{align*}
\end{proof}
The next example applies the preceding lemma to bound 
the $p$-Schatten-norm, when $\,p\,\ge\, 2\,$, by the spectral norm and the intrinsic dimension.
  \begin{example*}
    Let 
    $
    \,
      \mathbf{B} \in \mathbb{C}^{m\times n}
      \,
    $
    be any rectangular matrix and let $\,p\ge 2\,$.
    Then 
    $
    \,
      \varphi(x)
      \,
      :=
      \,
      \left| x \right|
      \,
    $
  defines a convex function with $\,\varphi(0)\,=\,0\,$.
  The intrinsic dimension lemma yields
  \begin{gather*}
    \norm{\mathbf{B}}^{\,p}_p
    \ 
    =
    \ 
    \mathrm{tr}
      \left|
      \mathbf{B}^* \mathbf{B}
      \right|
      ^{\,p/2}
      \ 
    \le
    \ 
    \mathrm{intdim}
    \,
      \mathbf{B}^* \mathbf{B}
      \ 
      \cdot
      \ 
      \norm{
      \mathbf{B}^* \mathbf{B}
    }^{\,p/2}
    \ 
    =
    \ 
    \mathrm{intdim}
    \,
      \mathbf{B}^* \mathbf{B}
      \ 
      \cdot
      \ 
      \norm{
        \mathbf{B}
    }^{\,p}
    \,.
  \end{gather*}
  If, additionally, $\mathbf{B}$ is self-adjoint and positive-semidefinite
  then it holds 
  \begin{gather*}
    \mathrm{tr}
    \,
      \mathbf{B}^* \mathbf{B}
      \ 
    =
    \ 
    \mathrm{tr}
    \,
    \mathbf{B}^{\,2}
    \ 
    =
    \ 
    \sum_{i=1}^{n} 
    \,
    \lambda_i^{\,2}
    \ 
    \le 
    \ 
    \left( 
      \,
    \sum_{i=1}^{n} 
    \,
    \lambda_i
    \,
    \right)
    ^{\!2}
    \ 
    =
    \ 
    \left( 
    \mathrm{tr}
    \,
    \mathbf{B}
    \right)
    ^{\,2}
    \,,
  \end{gather*}
  and consequently
  \begin{gather*}
    \norm{\mathbf{B}}^{\,p}_{p}
    \ 
    \le
    \ 
    \left( 
    \mathrm{intdim}
    \,
    \mathbf{B}
    \right)
    ^{\,2}
    \ 
      \cdot
      \ 
      \norm{
        \mathbf{B}
    }^{\,p}
    \,.
  \end{gather*}
  \end{example*}


  \todo[color=red!40, inline]{Apply to estimate \cite[above (A.4)]{Chen2012} to get intrinsic dimension version of Rosenthal-Pinelis inequality.}
\begin{takeaways}
  The notion of intrinsic dimension is useful when bounding convex
  functions of a positive-semidefinite matrix by its spectral norm.
  We saw how to derive bounds on the $p\,$-Schatten-norm when $\,p\,\ge\, 2$.
\end{takeaways}



\chapter{Empirical Processes}
\section{A Primer on Empirical Processes}
\todo[color=red!40,inline]{Add outer probability calculus. \cite{vaart2013} p.6}

Let 
$
(
\mathbb{D}
,
d
)
$
be a metric space, and let 
$
(\P_n)_{n\in \mathbb{N}}
  \P
$
be (Borel) probability measures
on
$
(
\mathbb{D}
,
\mathcal{D}
)
,
$
where 
$\mathcal{D}$
is the Borel $\sigma$-algebra on $\mathbb{D},$
the smallest $\sigma$-algebra
containing all open sets.
Then the sequence 
$\P_n$
\textbf{converges weakly}
to 
$\P,$
which we denote as $\P_n \rightsquigarrow \P,$
if and only if 
\begin{gather}
  \int_\mathbb{D}
  f
  \text{d}
  \P_n
  \to
  \int_\mathbb{D}
  f
  \text{d}
  \P
  \qquad
  \text{for all}
  \ 
  f
  \in
  C_b(\mathbb{D})
  .
\end{gather}
Here 
$
  C_b(\mathbb{D})
$
denotes the set of all bounded, continuous, real functions on $\mathbb{D}.$
Equivalently, if 
$X_n$ and $X$
are 
$\mathbb{D}$-valued
random variables with distribution 
$\P_n$ and $\P$
respectively, then 
$X_n \to X$
if and only if 
\begin{gather}
  \E[f(X_n)]
  \to
  \E[f(X)]
  \qquad
  \text{for all}
  \ 
  f
  \in
  C_b(\mathbb{D})
  .
\end{gather}

This definitions yield the classical theory of weak convergence.
For a modern treatment see \cite{Klenke2020}.

The classical theory requires that 
$\P_n$
is defined, for each $n\in \mathbb{N},$
on the Borel $\sigma$-algebra $\mathcal{D},$
or, equivalently, that $X_n$ is a Borel measurable map for each $n\in \mathbb{N}.$
If 
$
(
\Omega_n,
\mathcal{A}_n,
\P_n
)
$
are the underlying probability spaces on which the maps 
$X_n$
are defined, this means that
$X_n^{-1}(D)\in \mathcal{A}_n$
for every Borel set $D \in \mathcal{D}.$
This required measurability usually holds when $\mathbb{D}$
is a separable metric space such as $\R^k$
or $C([0,1])$ with the supremum metric.

However, this apparently modest requirement can and does easily fail when the metric space $\mathbb{D}$ is not separable.

\begin{example}
  \emph{\cite[Problem~1.7.3]{vaart2013}}
  Let $\mathbb{D}=D([0,1])$ be the \textbf{Skorohod space} of all right-continuous functions on $[0,1]$
  with left limits endowed with the metric induced by the supremum norm.
  Define 
  $
    X:
    [0,1]
    \to
    \mathbb{D}
    ,\ 
    \omega
    \mapsto
    \mathbf{1}_{[\omega,1]}
    .
  $
  If we equip $[0,1]$ with the Borel $\sigma$-algebra 
  $\mathcal{B}([0,1])$, then 
  $X$ is not measurable. To see this, let $B_s$ be the open ball of radius $1/2$ in $\mathbb{D}$ around the function $\mathbf{1}_{[s,1]}.$
  Now $X(\omega)\in B_s$
  if and only if $\omega=s.$ Indeed, if $\omega\neq s$ there exists an $x$ between $\omega$ and $s$ such that the difference of the indicator functions is 1 at $x$. Conversely, if the distance is greater than
  $1/2$ at a point $x\in [0,1]$, it is because $x$ lies between $\omega$ and $s$ and the indicator functions have difference 1.
  Since arbitrary (even uncountable) unions of open sets are open,
  we get for every $S\subseteq [0,1]$ the open set 
  $
  G
  :=
  \bigcup_{s\in S}
    B_s
    \in \mathcal{D}
    .
  $
  It follows
  $
  X^{-1}(G)=S 
  \ 
  \text{for all}
  \ 
  S \subseteq [0,1]
  .
  $
  Since not all subsets of $[0,1]$ are measurable, we have
  $
  X^{-1}(\mathcal{D})\nsubseteq \mathcal{B}([0,1])
.
$
But then $X$ is not measurable. The $\sigma$-algebra $\mathcal{D}$ is to large.
\end{example}



Let 
$
  \left( 
    \Omega,
    \mathcal{A},
    \P
  \right)
$
be a probability space
and
$
  \left( 
    \mathcal{X},
    \Sigma
  \right)
$
a measurable space.
Let 
$
  X_j
  :
  \left( 
    \Omega,
    \mathcal{A},
    \P
  \right)
  \to
  \left( 
    \mathcal{X},
    \Sigma
  \right)
  ,
  j=1,\ldots, n
$ be independent and identically-distributed (i.i.d.)
random variables
with probability distribution $\P_X$ 
and
$\mathcal{F}$ a family of measurable functions
$
  f:
  \left( 
    \mathcal{X},
    \Sigma
  \right)
    \to
  \left( 
    \R,
    \mathcal{B}(\R)
  \right)
$.
Consider the map
\begin{gather}
  f
  \mapsto
  G_n f
  :=
  \sqrt{n}
  \left( 
    \frac{1}{n}
    \sum_{i = 1}^{n}
      f(X_i)
    -
    \P_X f
  \right),
\end{gather}
where
$
  \P_X f 
  :=
  \int_\mathcal{X} f \text{d}\P_X.
$
We call 
$
  \left( 
    G_n f
  \right)_{f \in \mathcal{F}}
$
the empirical process indexed by $\mathcal{F}$.
Furthermore
\begin{gather}
  \norm{G_n f}_\mathcal{F}
  :=
  \sup_
        { f \in \mathcal{F}}
        \left|
          G_n f
        \right|
        .
\end{gather}




\section{Maximal Inequalities}
%%%%%%%%%%%%%%%%
% INTRODUCTION %
%%%%%%%%%%%%%%%%

Let 
$
  \left( 
    \Omega,
    \mathcal{A},
    \P
  \right)
$
be a probability space
and
$
  \left( 
    \mathcal{X},
    \Sigma
  \right)
$
a measurable space.
Let 
$
  X_j
  :
  \left( 
    \Omega,
    \mathcal{A},
    \P
  \right)
  \to
  \left( 
    \mathcal{X},
    \Sigma
  \right)
  ,
  j=1,\ldots, n
$ be independent and identically-distributed (i.i.d.)
random variables
with probability distribution $\P_X$ 
and
$\mathcal{F}$ a family of measurable functions
$
  f:
  \left( 
    \mathcal{X},
    \Sigma
  \right)
    \to
  \left( 
    \R,
    \mathcal{B}(\R)
  \right)
$.
Consider the map
\begin{gather}
  f
  \mapsto
  G_n f
  :=
  \sqrt{n}
  \left( 
    \frac{1}{n}
    \sum_{i = 1}^{n}
      f(X_i)
    -
    \P_X f
  \right),
\end{gather}
where
$
  \P_X f 
  :=
  \int_\mathcal{X} f \text{d}\P_X.
$
We call 
$
  \left( 
    G_n f
  \right)_{f \in \mathcal{F}}
$
the empirical process indexed by $\mathcal{F}$.
Furthermore
\begin{gather}
  \norm{G_n f}_\mathcal{F}
  :=
  \sup_
        { f \in \mathcal{F}}
        \left|
          G_n f
        \right|
        .
\end{gather}

\begin{lemma}
  \emph{(Bernstein Inequality for Empirical Processes)}
  For any bounded, measurable function $f$
  it holds for all $t > 0$
  \begin{gather}
    \P 
    \left(
      \left| 
        G_n f
      \right|
      >
      t
    \right)
    \le
    2
    \exp
    \left( 
      - \frac{1}{4}
      \frac{t^2}
      {
        \P_X(f^2)
        +
        t
        \norm{f}_\infty
        /
        \sqrt{n}
      }
    \right)
  \end{gather}
\end{lemma}
\begin{proof}
  By the Markov inequality it holds for all $\lambda > 0$
  \begin{gather}
    \P
    \left( 
      G_n f 
      > 
      t
    \right)
    \le
    e^{-\lambda t}
    \E
    \exp
    \left( 
      \lambda
      G_n f 
    \right)
  \end{gather}
\end{proof}


\begin{lemma}
  For any finite class $\mathcal{F}$ of bounded, measurable, square-integrable functions, with $\left| \mathcal{F} \right|$ elements, it holds
  \begin{gather}
    \E \norm{G_n f}_\mathcal{F}
    \lesssim
    \max_
          {f \in \mathcal{F}}
          \frac{\norm{f}_\infty}{\sqrt{n}}
          \log
          \left( 1 + \left| \mathcal{F} \right| \right)
          +
          \max_
              {f \in \mathcal{F}}
              \norm{f}_{\P,2}
              \sqrt{
                \log \left( 
                  1 + \left| \mathcal{F} \right|
                \right)
              }
              .
  \end{gather}
\end{lemma}

\begin{lemma}
  For any class 
  $\mathcal{F}$
  of measurable functions 
  $
    f:
    \mathcal{X}
    \to \R
  $
  such that 
  $
    \P f^2
    <
    \delta^2
  $
  for every $f,$
  we have,
  with
  $
    a(\delta)
    =
    \delta
    /
    \sqrt{
      \mathrm{Log}
    N
    _{[]}
    (
    \delta
    ,
    \mathcal{F}
    ,
    L_2(\P)
    )
    }
    ,
  $
  and $F$
  an envelope function,

  \begin{gather}
    \E^*
    _\P
    [
    \norm{
      \G
      _n
      }
      _\mathcal{F}
    ]
    \lesssim
    J
    _{[]}
    (
      \delta
    ,
    \mathcal{F}
    ,
    L_2(\P)
    )
    +
    \sqrt{n}
    \P^*
    F
    \left\{ 
      F > \sqrt{n}
      a(\delta)
    \right\}
    .
  \end{gather}
\end{lemma}


\begin{corollary}
  For any class $\mathcal{F}$ of measurable functions with envelope function $F,$
  \begin{gather}
    \E^*
    _\P
    [
    \norm{
      \G
      _n
      }
      _\mathcal{F}
    ]
    \lesssim
    J
    _{[]}
    (
    \norm{
      F
    }
    _{\P,2}
    ,
    \mathcal{F}
    ,
    L_2(\P)
    )
    .
  \end{gather}
\end{corollary}





\section{Functional Delta Method}
\begin{definition}
  A map 
  $
  \phi:
  \mathbb{D}_\phi
  \to 
  \mathbb{E}
  ,
  $
  defined on a subset 
  $
  \mathbb{D}_\phi
  $
  of a normed space
  $\mathbb{D}$
  that contains 
  $\theta,$
  is called 
  \textbf{Hadamard diffenertiable}
  at $\theta$
  if there exists a continuous,
  linear map
  $
  \phi_\theta^{'}
    :
    \mathbb{D}
    \to 
    \mathbb{E}
  $
  such that
  \begin{gather}
    \norm{
      \frac{
        \phi(\theta + t h_t)
        -
        \phi(\theta)
      }{
        t
      }
      -
      \phi^{'}_\theta
      (h)
    }_\mathbb{E}
    \to
    0
    \quad
    \text{as}
    \ 
    t\searrow 0
    \ 
    \text{for all}
    \ 
    h_t \to h
  \end{gather}
  $
    \text{such that $\theta + th_t$ is contained in $\mathbb{D}_\phi$ for all small $t>0.$}
  $
\end{definition}


\begin{ftheorem}
  \emph{(Delta Method)}
  Let 
  $
    \mathbb{D}
    \ \text{and}
    \ 
    \mathbb{E}
  $
  be normed linear spaces.
  Let
  $
    \phi
    :
    \mathbb{D}_\phi
    \subseteq
    \mathbb{D}
    \to
    \mathbb{E}
  $
  be Hadamard differentiable it $\theta$
  tangentially to 
  $\mathbb{D}_0.$
  Let
  $
    T_n
    :
    \Omega_n
    \to
    \mathbb{D}_\phi
  $
  be maps such that 
  $
    r_n
    (T_n - \theta)
    \rightsquigarrow
    T
  $
  for some sequence of numbers $r_n \to \infty$
  and a random element $T$
  that takes its values in $\mathbb{D}_0.$
  Then 
  $
    r_n(\phi(T_n)-\phi(\theta))
    \rightsquigarrow
    \phi^{'}
    _\theta
    (T)
    .
  $
  If 
  $
    \phi^{'}
    _\theta
  $
  is defined and continuous on the whole space $\mathbb{D},$
  then we also have 
  $
    r_n
    (
      \phi(T_n)
      -
      \phi(\theta)
    )
    =
    \phi^{'}
    _\theta
    (
    r_n
    (
    T_n
    -
    \theta
    )
    )
    +
    o_\P
    (1)
    .
  $
\end{ftheorem}
\begin{proof}
  \cite[Theorem~20.8]{Vaart1998}
\end{proof}


\chapter{Simple yet useful Calculations} 
\begin{proposition}
  Let 
  $f : \R^n \to \R$ 
  be continuous such that 
  a minimum $x^*$ exists and is unique.
  Then 
  for all $y \in \R^n$ and $C>0$ 
  it follows
    \begin{gather}
      \inf_{\norm{\Delta}=C} f(y+\Delta) - f(y) > 0 \qquad
      \Rightarrow \qquad 
      \norm{x^* - y} \le C.
    \end{gather}
\end{proposition}


\begin{proof}
Since 
$\mathcal{C}:=\left\{ \norm{\Delta}\le C \right\}$
is compact and
\begin{gather*}
  f(x^*) \le f(y) <  \inf_{\norm{\Delta}=C} f(y+\Delta)
\end{gather*}
the continious function $f(y+\,\cdot\,)$ has a minimum in 
$\overset{\circ}{\mathcal{C}}:=\left\{ \norm{\Delta} < C \right\}$. 
Since 
$x^*$ is the unique minimum of $f$
there exists $\Delta^* \in \overset{\circ}{\mathcal{C}}$ 
such that 
$x^* - y = \Delta^*$.
We conclude that
$\norm{x^* - y} \le C$.
\end{proof}

 %%%%%%%%%%%%%%%%%%%%%%%%%%%%%%%%%%%%%%%%%%%%%%%%%%%%%%%%%%%%%%%%


\begin{theorem}
  \emph{(Multivariate Taylor Theorem)}
  Let 
  $f \in C^2(\R^n, \R)$.
  Then 
  for all $x, \Delta \in \R^n$
  there exists $\xi \in [0,1]$
  such that 
  it holds
  \begin{gather}
    \label{syu_taylor}
    f(x + \Delta) - f(x)
    =
    \sum_{i = 1}^{n} \frac{\partial f(x)}{\partial x_i} \Delta_i    
    + \sum_{i,j = 1}^{n}  \frac{\partial^2 f(x + \xi \Delta)}{\partial x_i \partial x_j} \Delta_i\Delta_j
    + \frac{1}{2}\sum_{i = 1}^{n} \frac{\partial^2 f(x + \xi \Delta)}{\partial x_i^2}\Delta_i^2 
  \end{gather}
\end{theorem}


\begin{corollary}
  Let 
  $f\in C^2(\R)$. 
  Then
  for all $a,x,\Delta \in \R^n$ 
  there exist $\xi \in [0,1]$ 
  such that it holds
  \begin{gather}
    \label{syu_2.result}
    f(a^T (x + \Delta)) - f(a^T x) = 
    f^{'}(a^T x)\, \Delta^T a + 
    \frac{1}{2}f^{''}(a^T (x + \xi \Delta))\, \Delta^T A\ \Delta,
  \end{gather}
  where 
  $A:= a a^T \in \R^{n \times n}$ .
\end{corollary}

\begin{proof}
  By the chain rule 
  we have
  for all $a, x, \Delta \in \R^n$ and $\xi \in [0,1]$
  \begin{gather}
    \label{syu_2.0}
   \frac{\partial^2  f(a^T (x + \xi \Delta)) }{\partial x_i \partial x_j}=
    f^{''}(a^T (x + \xi \Delta))\, a_i a_j.
  \end{gather}
  Since 
  $A:= a a^T$ is symmetric
  we have
  \begin{gather}
    \label{syu_2.1}
    \Delta^T A\ \Delta 
    =
    2
    \sum_{
    \begin{smallmatrix}
      i,j = 1 \\ 
      i \neq j
    \end{smallmatrix}
    }
    ^{n}
    a_i a_j \Delta_i \Delta_j
    + 
    \sum_{i=1}^{n} 
    a_i^2 \Delta_i^2.
  \end{gather}
  Plugging \eqref{syu_2.0} and \eqref{syu_2.1} into \eqref{syu_taylor}
  yields
  \eqref{syu_2.result}.
\end{proof}




%makeindex main.nlo -s nomencl.ist -o main.nls

\nomenclature{$\P$}{generic probability measure}
\nomenclature{$\overset{\mathcal{D}}{\longrightarrow}$}{convergence of distributions}

\printnomenclature

\bibliography{literature}{}
\bibliographystyle{alpha}
\end{document}
